% Options for packages loaded elsewhere
\PassOptionsToPackage{unicode}{hyperref}
\PassOptionsToPackage{hyphens}{url}
\documentclass[
  11pt,
]{article}
\usepackage{xcolor}
\usepackage[margin=2.5cm]{geometry}
\usepackage{amsmath,amssymb}
\setcounter{secnumdepth}{-\maxdimen} % remove section numbering
\usepackage{iftex}
\ifPDFTeX
  \usepackage[T1]{fontenc}
  \usepackage[utf8]{inputenc}
  \usepackage{textcomp} % provide euro and other symbols
\else % if luatex or xetex
  \usepackage{unicode-math} % this also loads fontspec
  \defaultfontfeatures{Scale=MatchLowercase}
  \defaultfontfeatures[\rmfamily]{Ligatures=TeX,Scale=1}
\fi
\usepackage{lmodern}
\ifPDFTeX\else
  % xetex/luatex font selection
\fi
% Use upquote if available, for straight quotes in verbatim environments
\IfFileExists{upquote.sty}{\usepackage{upquote}}{}
\IfFileExists{microtype.sty}{% use microtype if available
  \usepackage[]{microtype}
  \UseMicrotypeSet[protrusion]{basicmath} % disable protrusion for tt fonts
}{}
\makeatletter
\@ifundefined{KOMAClassName}{% if non-KOMA class
  \IfFileExists{parskip.sty}{%
    \usepackage{parskip}
  }{% else
    \setlength{\parindent}{0pt}
    \setlength{\parskip}{6pt plus 2pt minus 1pt}}
}{% if KOMA class
  \KOMAoptions{parskip=half}}
\makeatother
\usepackage{color}
\usepackage{fancyvrb}
\newcommand{\VerbBar}{|}
\newcommand{\VERB}{\Verb[commandchars=\\\{\}]}
\DefineVerbatimEnvironment{Highlighting}{Verbatim}{commandchars=\\\{\}}
% Add ',fontsize=\small' for more characters per line
\usepackage{framed}
\definecolor{shadecolor}{RGB}{248,248,248}
\newenvironment{Shaded}{\begin{snugshade}}{\end{snugshade}}
\newcommand{\AlertTok}[1]{\textcolor[rgb]{0.94,0.16,0.16}{#1}}
\newcommand{\AnnotationTok}[1]{\textcolor[rgb]{0.56,0.35,0.01}{\textbf{\textit{#1}}}}
\newcommand{\AttributeTok}[1]{\textcolor[rgb]{0.13,0.29,0.53}{#1}}
\newcommand{\BaseNTok}[1]{\textcolor[rgb]{0.00,0.00,0.81}{#1}}
\newcommand{\BuiltInTok}[1]{#1}
\newcommand{\CharTok}[1]{\textcolor[rgb]{0.31,0.60,0.02}{#1}}
\newcommand{\CommentTok}[1]{\textcolor[rgb]{0.56,0.35,0.01}{\textit{#1}}}
\newcommand{\CommentVarTok}[1]{\textcolor[rgb]{0.56,0.35,0.01}{\textbf{\textit{#1}}}}
\newcommand{\ConstantTok}[1]{\textcolor[rgb]{0.56,0.35,0.01}{#1}}
\newcommand{\ControlFlowTok}[1]{\textcolor[rgb]{0.13,0.29,0.53}{\textbf{#1}}}
\newcommand{\DataTypeTok}[1]{\textcolor[rgb]{0.13,0.29,0.53}{#1}}
\newcommand{\DecValTok}[1]{\textcolor[rgb]{0.00,0.00,0.81}{#1}}
\newcommand{\DocumentationTok}[1]{\textcolor[rgb]{0.56,0.35,0.01}{\textbf{\textit{#1}}}}
\newcommand{\ErrorTok}[1]{\textcolor[rgb]{0.64,0.00,0.00}{\textbf{#1}}}
\newcommand{\ExtensionTok}[1]{#1}
\newcommand{\FloatTok}[1]{\textcolor[rgb]{0.00,0.00,0.81}{#1}}
\newcommand{\FunctionTok}[1]{\textcolor[rgb]{0.13,0.29,0.53}{\textbf{#1}}}
\newcommand{\ImportTok}[1]{#1}
\newcommand{\InformationTok}[1]{\textcolor[rgb]{0.56,0.35,0.01}{\textbf{\textit{#1}}}}
\newcommand{\KeywordTok}[1]{\textcolor[rgb]{0.13,0.29,0.53}{\textbf{#1}}}
\newcommand{\NormalTok}[1]{#1}
\newcommand{\OperatorTok}[1]{\textcolor[rgb]{0.81,0.36,0.00}{\textbf{#1}}}
\newcommand{\OtherTok}[1]{\textcolor[rgb]{0.56,0.35,0.01}{#1}}
\newcommand{\PreprocessorTok}[1]{\textcolor[rgb]{0.56,0.35,0.01}{\textit{#1}}}
\newcommand{\RegionMarkerTok}[1]{#1}
\newcommand{\SpecialCharTok}[1]{\textcolor[rgb]{0.81,0.36,0.00}{\textbf{#1}}}
\newcommand{\SpecialStringTok}[1]{\textcolor[rgb]{0.31,0.60,0.02}{#1}}
\newcommand{\StringTok}[1]{\textcolor[rgb]{0.31,0.60,0.02}{#1}}
\newcommand{\VariableTok}[1]{\textcolor[rgb]{0.00,0.00,0.00}{#1}}
\newcommand{\VerbatimStringTok}[1]{\textcolor[rgb]{0.31,0.60,0.02}{#1}}
\newcommand{\WarningTok}[1]{\textcolor[rgb]{0.56,0.35,0.01}{\textbf{\textit{#1}}}}
\usepackage{longtable,booktabs,array}
\newcounter{none} % for unnumbered tables
\usepackage{calc} % for calculating minipage widths
% Correct order of tables after \paragraph or \subparagraph
\usepackage{etoolbox}
\makeatletter
\patchcmd\longtable{\par}{\if@noskipsec\mbox{}\fi\par}{}{}
\makeatother
% Allow footnotes in longtable head/foot
\IfFileExists{footnotehyper.sty}{\usepackage{footnotehyper}}{\usepackage{footnote}}
\makesavenoteenv{longtable}
\usepackage{graphicx}
\makeatletter
\newsavebox\pandoc@box
\newcommand*\pandocbounded[1]{% scales image to fit in text height/width
  \sbox\pandoc@box{#1}%
  \Gscale@div\@tempa{\textheight}{\dimexpr\ht\pandoc@box+\dp\pandoc@box\relax}%
  \Gscale@div\@tempb{\linewidth}{\wd\pandoc@box}%
  \ifdim\@tempb\p@<\@tempa\p@\let\@tempa\@tempb\fi% select the smaller of both
  \ifdim\@tempa\p@<\p@\scalebox{\@tempa}{\usebox\pandoc@box}%
  \else\usebox{\pandoc@box}%
  \fi%
}
% Set default figure placement to htbp
\def\fps@figure{htbp}
\makeatother
\setlength{\emergencystretch}{3em} % prevent overfull lines
\providecommand{\tightlist}{%
  \setlength{\itemsep}{0pt}\setlength{\parskip}{0pt}}
\usepackage{bookmark}
\IfFileExists{xurl.sty}{\usepackage{xurl}}{} % add URL line breaks if available
\urlstyle{same}
\hypersetup{
  pdftitle={Causal Discovery pada Dataset Pima Indians Diabetes: Implementasi Algoritma PC-GES},
  pdfauthor={Muhammad Dhiauddin - 25917024},
  hidelinks,
  pdfcreator={LaTeX via pandoc}}

\title{Causal Discovery pada Dataset Pima Indians Diabetes: Implementasi
Algoritma PC-GES}
\author{Muhammad Dhiauddin - 25917024}
\date{22 June 2026}

\begin{document}
\maketitle

{
\setcounter{tocdepth}{3}
\tableofcontents
}
\begin{Shaded}
\begin{Highlighting}[]
\CommentTok{\# ── Jalankan blok ini SATU KALI untuk menginstal semua dependensi ──────────────}
\FunctionTok{install.packages}\NormalTok{(}\FunctionTok{c}\NormalTok{(}\StringTok{"bnlearn"}\NormalTok{, }\StringTok{"igraph"}\NormalTok{, }\StringTok{"dagitty"}\NormalTok{,}
                   \StringTok{"ggplot2"}\NormalTok{, }\StringTok{"dplyr"}\NormalTok{, }\StringTok{"corrplot"}\NormalTok{, }\StringTok{"knitr"}\NormalTok{))}

\CommentTok{\# Rgraphviz (opsional, alternatif visualisasi via Bioconductor)}
\ControlFlowTok{if}\NormalTok{ (}\SpecialCharTok{!}\FunctionTok{requireNamespace}\NormalTok{(}\StringTok{"BiocManager"}\NormalTok{, }\AttributeTok{quietly =} \ConstantTok{TRUE}\NormalTok{))}
  \FunctionTok{install.packages}\NormalTok{(}\StringTok{"BiocManager"}\NormalTok{)}
\NormalTok{BiocManager}\SpecialCharTok{::}\FunctionTok{install}\NormalTok{(}\FunctionTok{c}\NormalTok{(}\StringTok{"Rgraphviz"}\NormalTok{, }\StringTok{"graph"}\NormalTok{))}
\end{Highlighting}
\end{Shaded}

\begin{Shaded}
\begin{Highlighting}[]
\FunctionTok{library}\NormalTok{(bnlearn)}
\FunctionTok{library}\NormalTok{(igraph)}
\FunctionTok{library}\NormalTok{(dagitty)}
\FunctionTok{library}\NormalTok{(ggplot2)}
\FunctionTok{library}\NormalTok{(dplyr)}
\FunctionTok{library}\NormalTok{(corrplot)}
\FunctionTok{library}\NormalTok{(knitr)}
\FunctionTok{set.seed}\NormalTok{(}\DecValTok{42}\NormalTok{)}
\end{Highlighting}
\end{Shaded}

\begin{center}\rule{0.5\linewidth}{0.5pt}\end{center}

\textbf{Identitas Mahasiswa:}

\begin{itemize}
\tightlist
\item
  \textbf{Nama:} Muhammad Dhiauddin
\item
  \textbf{NIM:} 25917024
\item
  \textbf{Konsentrasi:} Sains Data - Profesional
\item
  \textbf{Github:}
  \url{https://github.com/Mudhya19/Project-Causal-Discovery-PC-Algorithm-vs-GES.git}
\item
  \textbf{Kaggle:}
  \url{https://www.kaggle.com/datasets/uciml/pima-indians-diabetes-database}
\end{itemize}

\begin{center}\rule{0.5\linewidth}{0.5pt}\end{center}

\section{Pendahuluan}\label{pendahuluan}

Diabetes melitus tipe 2 merupakan salah satu penyakit metabolik kronis
dengan prevalensi tertinggi di dunia. World Health Organization (WHO)
mencatat lebih dari 422 juta orang menderita diabetes, dengan sebagian
besar beban penyakit berada di negara berkembang. Kemampuan untuk
memahami \textbf{faktor kausal} - bukan sekadar faktor yang berkorelasi
- sangat penting agar intervensi kesehatan dapat dirancang secara lebih
efektif. Sebagai contoh, mengetahui bahwa penurunan BMI secara
\emph{kausal} menurunkan risiko diabetes memberikan justifikasi yang
lebih kuat untuk program penurunan berat badan dibandingkan sekadar
mengetahui bahwa BMI dan diabetes berkorelasi positif.

Pendekatan \emph{causal discovery} memungkinkan kita memperoleh struktur
kausal secara data-driven dari data observasional, tanpa memerlukan
eksperimen terkontrol (RCT). Salah satu algoritma yang paling banyak
digunakan adalah \textbf{PC Algorithm} (Spirtes et al., 2000), yang
bekerja dengan menjalankan serangkaian uji independensi kondisional (CI
test) untuk membangun \emph{Completed Partially Directed Acyclic Graph}
(CPDAG) - representasi kelas ekivalensi Markov dari semua DAG yang
kompatibel dengan data.

\textbf{Pertanyaan kausal utama} yang dirumuskan dalam studi ini adalah:

\begin{quote}
\emph{``Apakah Indeks Massa Tubuh (BMI) dan kadar glukosa darah secara
kausal mempengaruhi risiko diabetes (Outcome), atau hubungan tersebut
dijembatani atau dikonfound oleh variabel lain seperti usia (Age) dan
riwayat kehamilan (Pregnancies)?''}
\end{quote}

Dataset yang digunakan adalah \textbf{Pima Indians Diabetes Database}
yang bersumber dari National Institute of Diabetes and Digestive and
Kidney Diseases (NIDDK), tersedia secara publik di Kaggle
(\url{https://www.kaggle.com/datasets/uciml/pima-indians-diabetes-database}).
Dataset ini mencatat pengukuran medis dari 768 wanita keturunan suku
Pima Indian berusia minimal 21 tahun.

\begin{center}\rule{0.5\linewidth}{0.5pt}\end{center}

\section{Metode}\label{metode}

\subsection{Deskripsi Variabel}\label{deskripsi-variabel}

Dataset terdiri atas 9 variabel dengan tipe data yang dijelaskan pada
Tabel 1.

\begin{longtable}[]{@{}
  >{\raggedleft\arraybackslash}p{(\linewidth - 6\tabcolsep) * \real{0.0312}}
  >{\raggedright\arraybackslash}p{(\linewidth - 6\tabcolsep) * \real{0.2604}}
  >{\raggedright\arraybackslash}p{(\linewidth - 6\tabcolsep) * \real{0.1458}}
  >{\raggedright\arraybackslash}p{(\linewidth - 6\tabcolsep) * \real{0.5625}}@{}}
\caption{Deskripsi Variabel Dataset Pima Indians
Diabetes}\tabularnewline
\toprule\noalign{}
\begin{minipage}[b]{\linewidth}\raggedleft
No
\end{minipage} & \begin{minipage}[b]{\linewidth}\raggedright
Variabel
\end{minipage} & \begin{minipage}[b]{\linewidth}\raggedright
Tipe
\end{minipage} & \begin{minipage}[b]{\linewidth}\raggedright
Deskripsi
\end{minipage} \\
\midrule\noalign{}
\endfirsthead
\toprule\noalign{}
\begin{minipage}[b]{\linewidth}\raggedleft
No
\end{minipage} & \begin{minipage}[b]{\linewidth}\raggedright
Variabel
\end{minipage} & \begin{minipage}[b]{\linewidth}\raggedright
Tipe
\end{minipage} & \begin{minipage}[b]{\linewidth}\raggedright
Deskripsi
\end{minipage} \\
\midrule\noalign{}
\endhead
\bottomrule\noalign{}
\endlastfoot
1 & Pregnancies & Kontinu & Jumlah kehamilan \\
2 & Glucose & Kontinu & Glukosa plasma 2 jam pasca tes toleransi oral
(mg/dL) \\
3 & BloodPressure & Kontinu & Tekanan darah diastolik (mm Hg) \\
4 & SkinThickness & Kontinu & Ketebalan lipatan kulit trisep (mm) \\
5 & Insulin & Kontinu & Kadar insulin serum 2 jam (µU/mL) \\
6 & BMI & Kontinu & Indeks Massa Tubuh (kg/m²) \\
7 & DiabetesPedigreeFunction & Kontinu & Skor riwayat diabetes
keluarga \\
8 & Age & Kontinu & Usia (tahun) \\
9 & Outcome & Biner Diskrit & Diagnosis diabetes (1 = Positif, 0 =
Negatif) \\
\end{longtable}

\subsection{Eksplorasi dan Persiapan
Data}\label{eksplorasi-dan-persiapan-data}

\subsubsection{Pembacaan Data dan Statistik Deskriptif
Awal}\label{pembacaan-data-dan-statistik-deskriptif-awal}

\begin{Shaded}
\begin{Highlighting}[]
\CommentTok{\# Pastikan file diabetes.csv berada di direktori yang tepat}
\NormalTok{diabetes\_raw }\OtherTok{\textless{}{-}} \FunctionTok{read.csv}\NormalTok{(}\StringTok{"dataset/diabetes.csv"}\NormalTok{, }\AttributeTok{stringsAsFactors =} \ConstantTok{FALSE}\NormalTok{)}

\CommentTok{\# Konversi semua kolom integer menjadi numeric untuk menghindari error di bnlearn}
\NormalTok{diabetes\_raw[] }\OtherTok{\textless{}{-}} \FunctionTok{lapply}\NormalTok{(diabetes\_raw, }\ControlFlowTok{function}\NormalTok{(x) }\ControlFlowTok{if}\NormalTok{(}\FunctionTok{is.integer}\NormalTok{(x)) }\FunctionTok{as.numeric}\NormalTok{(x) }\ControlFlowTok{else}\NormalTok{ x)}

\FunctionTok{cat}\NormalTok{(}\StringTok{"Dimensi dataset:"}\NormalTok{, }\FunctionTok{nrow}\NormalTok{(diabetes\_raw), }\StringTok{"baris x"}\NormalTok{,}
    \FunctionTok{ncol}\NormalTok{(diabetes\_raw), }\StringTok{"kolom}\SpecialCharTok{\textbackslash{}n}\StringTok{"}\NormalTok{)}
\end{Highlighting}
\end{Shaded}

\begin{verbatim}
## Dimensi dataset: 768 baris x 9 kolom
\end{verbatim}

\begin{Shaded}
\begin{Highlighting}[]
\CommentTok{\# Menampilkan tipe data dalam format tabel yang rapi}
\NormalTok{tipe\_data\_df }\OtherTok{\textless{}{-}} \FunctionTok{data.frame}\NormalTok{(}
  \AttributeTok{Variabel =} \FunctionTok{names}\NormalTok{(diabetes\_raw),}
  \AttributeTok{Tipe\_Data =} \FunctionTok{sapply}\NormalTok{(diabetes\_raw, class),}
  \AttributeTok{row.names =} \ConstantTok{NULL}
\NormalTok{)}
\NormalTok{knitr}\SpecialCharTok{::}\FunctionTok{kable}\NormalTok{(tipe\_data\_df, }\AttributeTok{caption =} \StringTok{"Tipe Data Setiap Kolom"}\NormalTok{, }\AttributeTok{row.names =} \ConstantTok{FALSE}\NormalTok{)}
\end{Highlighting}
\end{Shaded}

\begin{longtable}[]{@{}ll@{}}
\caption{Tipe Data Setiap Kolom}\tabularnewline
\toprule\noalign{}
Variabel & Tipe\_Data \\
\midrule\noalign{}
\endfirsthead
\toprule\noalign{}
Variabel & Tipe\_Data \\
\midrule\noalign{}
\endhead
\bottomrule\noalign{}
\endlastfoot
Pregnancies & numeric \\
Glucose & numeric \\
BloodPressure & numeric \\
SkinThickness & numeric \\
Insulin & numeric \\
BMI & numeric \\
DiabetesPedigreeFunction & numeric \\
Age & numeric \\
Outcome & numeric \\
\end{longtable}

\begin{Shaded}
\begin{Highlighting}[]
\CommentTok{\# Transpose tabel summary agar variabel menjadi baris, bukan kolom (menghindari kepanjangan ke kanan)}
\NormalTok{desc\_stats }\OtherTok{\textless{}{-}} \FunctionTok{t}\NormalTok{(}\FunctionTok{round}\NormalTok{(}\FunctionTok{do.call}\NormalTok{(cbind, }\FunctionTok{lapply}\NormalTok{(diabetes\_raw, summary)), }\DecValTok{2}\NormalTok{))}
\NormalTok{knitr}\SpecialCharTok{::}\FunctionTok{kable}\NormalTok{(desc\_stats, }\AttributeTok{caption =} \StringTok{"Statistik Deskriptif Awal (Sebelum Preprocessing)"}\NormalTok{, }\AttributeTok{row.names =} \ConstantTok{TRUE}\NormalTok{)}
\end{Highlighting}
\end{Shaded}

\begin{longtable}[]{@{}
  >{\raggedright\arraybackslash}p{(\linewidth - 12\tabcolsep) * \real{0.3676}}
  >{\raggedleft\arraybackslash}p{(\linewidth - 12\tabcolsep) * \real{0.0882}}
  >{\raggedleft\arraybackslash}p{(\linewidth - 12\tabcolsep) * \real{0.1176}}
  >{\raggedleft\arraybackslash}p{(\linewidth - 12\tabcolsep) * \real{0.1029}}
  >{\raggedleft\arraybackslash}p{(\linewidth - 12\tabcolsep) * \real{0.1029}}
  >{\raggedleft\arraybackslash}p{(\linewidth - 12\tabcolsep) * \real{0.1176}}
  >{\raggedleft\arraybackslash}p{(\linewidth - 12\tabcolsep) * \real{0.1029}}@{}}
\caption{Statistik Deskriptif Awal (Sebelum
Preprocessing)}\tabularnewline
\toprule\noalign{}
\begin{minipage}[b]{\linewidth}\raggedright
\end{minipage} & \begin{minipage}[b]{\linewidth}\raggedleft
Min.
\end{minipage} & \begin{minipage}[b]{\linewidth}\raggedleft
1st Qu.
\end{minipage} & \begin{minipage}[b]{\linewidth}\raggedleft
Median
\end{minipage} & \begin{minipage}[b]{\linewidth}\raggedleft
Mean
\end{minipage} & \begin{minipage}[b]{\linewidth}\raggedleft
3rd Qu.
\end{minipage} & \begin{minipage}[b]{\linewidth}\raggedleft
Max.
\end{minipage} \\
\midrule\noalign{}
\endfirsthead
\toprule\noalign{}
\begin{minipage}[b]{\linewidth}\raggedright
\end{minipage} & \begin{minipage}[b]{\linewidth}\raggedleft
Min.
\end{minipage} & \begin{minipage}[b]{\linewidth}\raggedleft
1st Qu.
\end{minipage} & \begin{minipage}[b]{\linewidth}\raggedleft
Median
\end{minipage} & \begin{minipage}[b]{\linewidth}\raggedleft
Mean
\end{minipage} & \begin{minipage}[b]{\linewidth}\raggedleft
3rd Qu.
\end{minipage} & \begin{minipage}[b]{\linewidth}\raggedleft
Max.
\end{minipage} \\
\midrule\noalign{}
\endhead
\bottomrule\noalign{}
\endlastfoot
Pregnancies & 0.00 & 1.00 & 3.00 & 3.85 & 6.00 & 17.00 \\
Glucose & 0.00 & 99.00 & 117.00 & 120.89 & 140.25 & 199.00 \\
BloodPressure & 0.00 & 62.00 & 72.00 & 69.11 & 80.00 & 122.00 \\
SkinThickness & 0.00 & 0.00 & 23.00 & 20.54 & 32.00 & 99.00 \\
Insulin & 0.00 & 0.00 & 30.50 & 79.80 & 127.25 & 846.00 \\
BMI & 0.00 & 27.30 & 32.00 & 31.99 & 36.60 & 67.10 \\
DiabetesPedigreeFunction & 0.08 & 0.24 & 0.37 & 0.47 & 0.63 & 2.42 \\
Age & 21.00 & 24.00 & 29.00 & 33.24 & 41.00 & 81.00 \\
Outcome & 0.00 & 0.00 & 0.00 & 0.35 & 1.00 & 1.00 \\
\end{longtable}

\begin{Shaded}
\begin{Highlighting}[]
\FunctionTok{par}\NormalTok{(}\AttributeTok{mfrow =} \FunctionTok{c}\NormalTok{(}\DecValTok{3}\NormalTok{, }\DecValTok{3}\NormalTok{), }\AttributeTok{mar =} \FunctionTok{c}\NormalTok{(}\DecValTok{3}\NormalTok{, }\DecValTok{3}\NormalTok{, }\DecValTok{2}\NormalTok{, }\DecValTok{1}\NormalTok{))}
\NormalTok{num\_cols }\OtherTok{\textless{}{-}} \FunctionTok{names}\NormalTok{(diabetes\_raw)}
\ControlFlowTok{for}\NormalTok{ (col }\ControlFlowTok{in}\NormalTok{ num\_cols) \{}
  \ControlFlowTok{if}\NormalTok{ (col }\SpecialCharTok{==} \StringTok{"Outcome"}\NormalTok{) \{}
    \FunctionTok{barplot}\NormalTok{(}\FunctionTok{table}\NormalTok{(diabetes\_raw[[col]]),}
            \AttributeTok{main =}\NormalTok{ col, }\AttributeTok{col =} \FunctionTok{c}\NormalTok{(}\StringTok{"\#AED6F1"}\NormalTok{, }\StringTok{"\#E74C3C"}\NormalTok{),}
            \AttributeTok{names.arg =} \FunctionTok{c}\NormalTok{(}\StringTok{"Negatif (0)"}\NormalTok{, }\StringTok{"Positif (1)"}\NormalTok{), }\AttributeTok{cex.main =} \FloatTok{0.9}\NormalTok{)}
\NormalTok{  \} }\ControlFlowTok{else}\NormalTok{ \{}
    \FunctionTok{hist}\NormalTok{(diabetes\_raw[[col]], }\AttributeTok{main =}\NormalTok{ col, }\AttributeTok{col =} \StringTok{"\#AED6F1"}\NormalTok{,}
         \AttributeTok{border =} \StringTok{"white"}\NormalTok{, }\AttributeTok{xlab =} \StringTok{""}\NormalTok{, }\AttributeTok{cex.main =} \FloatTok{0.9}\NormalTok{)}
\NormalTok{  \}}
\NormalTok{\}}
\end{Highlighting}
\end{Shaded}

\begin{figure}

{\centering \includegraphics[width=0.85\linewidth]{pc_algorithm_vs_ges_files/figure-latex/dist-plot-1} 

}

\caption{Distribusi Seluruh Variabel Sebelum Preprocessing}\label{fig:dist-plot}
\end{figure}

\begin{Shaded}
\begin{Highlighting}[]
\FunctionTok{par}\NormalTok{(}\AttributeTok{mfrow =} \FunctionTok{c}\NormalTok{(}\DecValTok{1}\NormalTok{, }\DecValTok{1}\NormalTok{))}
\end{Highlighting}
\end{Shaded}

\subsubsection{Penanganan Missing
Values}\label{penanganan-missing-values}

Lima variabel mengandung nilai \texttt{0} yang secara fisiologis tidak
mungkin terjadi: \textbf{Glucose} (glukosa 0 tidak mungkin pada individu
hidup), \textbf{BloodPressure} (tekanan darah diastolik 0 berarti
kematian klinis), \textbf{SkinThickness} (ketebalan kulit trisep 0 tidak
mungkin), \textbf{Insulin} (kadar insulin 0 sangat tidak realistis
kecuali dalam kondisi patologis ekstrem), dan \textbf{BMI} (BMI 0 tidak
mungkin pada manusia). Nilai-nilai ini merupakan \emph{missing values}
yang dikodekan sebagai \texttt{0}, kemungkinan akibat keterbatasan
sistem pencatatan medis pada era pengumpulan data.

\textbf{Justifikasi Strategi Imputasi - Median Imputation:} Dipilih
\emph{median imputation} karena (1) median lebih robust terhadap
distribusi skewed dan outlier dibandingkan mean; (2) variabel medis
seperti Insulin dan SkinThickness memiliki distribusi yang right-skewed
sehingga mean akan overestimate nilai sentral; (3) sederhana, dapat
direproduksi, dan tidak menambah variansi buatan. Sebagai alternatif,
\emph{Multiple Imputation by Chained Equations} (MICE) dapat
dipertimbangkan untuk memperoleh estimasi yang lebih tidak bias, namun
untuk keperluan algoritma causal discovery dengan fokus pada struktur
kausal, median imputation sudah memadai dan umum digunakan dalam
literatur (Spirtes et al., 2000).

\begin{Shaded}
\begin{Highlighting}[]
\NormalTok{zero\_cols }\OtherTok{\textless{}{-}} \FunctionTok{c}\NormalTok{(}\StringTok{"Glucose"}\NormalTok{, }\StringTok{"BloodPressure"}\NormalTok{, }\StringTok{"SkinThickness"}\NormalTok{, }\StringTok{"Insulin"}\NormalTok{, }\StringTok{"BMI"}\NormalTok{)}

\CommentTok{\# Salin dataset dan ganti 0 dengan NA}
\NormalTok{diabetes\_na }\OtherTok{\textless{}{-}}\NormalTok{ diabetes\_raw}
\ControlFlowTok{for}\NormalTok{ (col }\ControlFlowTok{in}\NormalTok{ zero\_cols) \{}
\NormalTok{  diabetes\_na[[col]][diabetes\_na[[col]] }\SpecialCharTok{==} \DecValTok{0}\NormalTok{] }\OtherTok{\textless{}{-}} \ConstantTok{NA}
\NormalTok{\}}

\CommentTok{\# Tampilkan jumlah missing values}
\NormalTok{mv\_df }\OtherTok{\textless{}{-}} \FunctionTok{data.frame}\NormalTok{(}
  \AttributeTok{Variabel   =} \FunctionTok{names}\NormalTok{(diabetes\_na),}
  \AttributeTok{Jumlah\_NA  =} \FunctionTok{colSums}\NormalTok{(}\FunctionTok{is.na}\NormalTok{(diabetes\_na)),}
  \AttributeTok{Persen\_NA  =} \FunctionTok{round}\NormalTok{(}\FunctionTok{colSums}\NormalTok{(}\FunctionTok{is.na}\NormalTok{(diabetes\_na)) }\SpecialCharTok{/} \FunctionTok{nrow}\NormalTok{(diabetes\_na) }\SpecialCharTok{*} \DecValTok{100}\NormalTok{, }\DecValTok{1}\NormalTok{)}
\NormalTok{)}
\FunctionTok{kable}\NormalTok{(mv\_df, }\AttributeTok{caption =} \StringTok{"Jumlah dan Persentase Missing Values per Variabel"}\NormalTok{,}
      \AttributeTok{col.names =} \FunctionTok{c}\NormalTok{(}\StringTok{"Variabel"}\NormalTok{, }\StringTok{"Jumlah NA"}\NormalTok{, }\StringTok{"Persen (\%)"}\NormalTok{),}
      \AttributeTok{row.names =} \ConstantTok{FALSE}\NormalTok{)}
\end{Highlighting}
\end{Shaded}

\begin{longtable}[]{@{}lrr@{}}
\caption{Jumlah dan Persentase Missing Values per
Variabel}\tabularnewline
\toprule\noalign{}
Variabel & Jumlah NA & Persen (\%) \\
\midrule\noalign{}
\endfirsthead
\toprule\noalign{}
Variabel & Jumlah NA & Persen (\%) \\
\midrule\noalign{}
\endhead
\bottomrule\noalign{}
\endlastfoot
Pregnancies & 0 & 0.0 \\
Glucose & 5 & 0.7 \\
BloodPressure & 35 & 4.6 \\
SkinThickness & 227 & 29.6 \\
Insulin & 374 & 48.7 \\
BMI & 11 & 1.4 \\
DiabetesPedigreeFunction & 0 & 0.0 \\
Age & 0 & 0.0 \\
Outcome & 0 & 0.0 \\
\end{longtable}

\begin{Shaded}
\begin{Highlighting}[]
\CommentTok{\# Lakukan Median Imputation}
\NormalTok{diabetes\_imp }\OtherTok{\textless{}{-}}\NormalTok{ diabetes\_na}
\NormalTok{medians\_used }\OtherTok{\textless{}{-}} \FunctionTok{numeric}\NormalTok{(}\FunctionTok{length}\NormalTok{(zero\_cols))}
\FunctionTok{names}\NormalTok{(medians\_used) }\OtherTok{\textless{}{-}}\NormalTok{ zero\_cols}

\ControlFlowTok{for}\NormalTok{ (col }\ControlFlowTok{in}\NormalTok{ zero\_cols) \{}
\NormalTok{  med\_val }\OtherTok{\textless{}{-}} \FunctionTok{median}\NormalTok{(diabetes\_imp[[col]], }\AttributeTok{na.rm =} \ConstantTok{TRUE}\NormalTok{)}
\NormalTok{  medians\_used[col] }\OtherTok{\textless{}{-}}\NormalTok{ med\_val}
\NormalTok{  diabetes\_imp[[col]][}\FunctionTok{is.na}\NormalTok{(diabetes\_imp[[col]])] }\OtherTok{\textless{}{-}}\NormalTok{ med\_val}
\NormalTok{\}}

\FunctionTok{cat}\NormalTok{(}\StringTok{"Nilai Median yang Digunakan untuk Imputasi:}\SpecialCharTok{\textbackslash{}n}\StringTok{"}\NormalTok{)}
\end{Highlighting}
\end{Shaded}

\begin{verbatim}
## Nilai Median yang Digunakan untuk Imputasi:
\end{verbatim}

\begin{Shaded}
\begin{Highlighting}[]
\NormalTok{medians\_df }\OtherTok{\textless{}{-}} \FunctionTok{data.frame}\NormalTok{(}\AttributeTok{Variabel =} \FunctionTok{names}\NormalTok{(medians\_used), }\AttributeTok{Median =}\NormalTok{ medians\_used, }\AttributeTok{row.names =} \ConstantTok{NULL}\NormalTok{)}
\NormalTok{knitr}\SpecialCharTok{::}\FunctionTok{kable}\NormalTok{(medians\_df, }\AttributeTok{caption =} \StringTok{"Nilai Median yang Digunakan untuk Imputasi"}\NormalTok{)}
\end{Highlighting}
\end{Shaded}

\begin{longtable}[]{@{}lr@{}}
\caption{Nilai Median yang Digunakan untuk Imputasi}\tabularnewline
\toprule\noalign{}
Variabel & Median \\
\midrule\noalign{}
\endfirsthead
\toprule\noalign{}
Variabel & Median \\
\midrule\noalign{}
\endhead
\bottomrule\noalign{}
\endlastfoot
Glucose & 117.0 \\
BloodPressure & 72.0 \\
SkinThickness & 29.0 \\
Insulin & 125.0 \\
BMI & 32.3 \\
\end{longtable}

\begin{Shaded}
\begin{Highlighting}[]
\FunctionTok{cat}\NormalTok{(}\StringTok{"}\SpecialCharTok{\textbackslash{}n}\StringTok{Verifikasi {-} Total NA setelah imputasi:"}\NormalTok{,}
    \FunctionTok{sum}\NormalTok{(}\FunctionTok{is.na}\NormalTok{(diabetes\_imp)), }\StringTok{"}\SpecialCharTok{\textbackslash{}n}\StringTok{"}\NormalTok{)}
\end{Highlighting}
\end{Shaded}

\begin{verbatim}
## 
## Verifikasi - Total NA setelah imputasi: 0
\end{verbatim}

\subsubsection{Penanganan Mixed Data:
Diskretisasi}\label{penanganan-mixed-data-diskretisasi}

\textbf{Justifikasi Pemilihan CI Test - G-Squared (Mutual Information)
Test:} Dataset ini merupakan \emph{mixed-type data} (8 variabel kontinu
+ 1 variabel biner diskrit). Terdapat tiga strategi penanganan:

\begin{enumerate}
\def\labelenumi{\arabic{enumi}.}
\tightlist
\item
  \textbf{Diskretisasi} semua variabel kontinu -\textgreater{} gunakan
  \emph{G-squared test}
\item
  \textbf{Kernel-based CI test (KCI)} - bekerja langsung pada data
  mixed, namun secara komputasional sangat mahal untuk 768 observasi
\item
  \textbf{Fisher's Z-test} - hanya tepat untuk data kontinu
  berdistribusi normal (Gaussian)
\end{enumerate}

Strategi yang dipilih adalah \textbf{diskretisasi} menggunakan fungsi
\texttt{discretize()} dari paket \texttt{bnlearn} dengan metode
\emph{quantile} (3 bin). Metode quantile dipilih karena: (a) memastikan
distribusi observasi yang relatif seimbang di setiap bin, mengurangi
risiko \emph{sparse cells} yang melemahkan uji statistik; (b) tidak
mengasumsikan distribusi tertentu sehingga cocok untuk variabel medis
yang umumnya skewed.

\textbf{Nama CI Test:} \emph{G-Squared Test} (ekivalen dengan
\emph{Mutual Information Test}), yaitu uji statistik yang mengikuti
distribusi \(\chi^2\) dan mengukur apakah dua variabel diskrit saling
bebas diberikan sekumpulan variabel kondisional. Statistik ujinya
adalah:
\[G^2 = 2 \sum_{i,j,k} n_{ijk} \ln\!\left(\frac{n_{ijk} \cdot n_{\cdot\cdot k}}{n_{i\cdot k} \cdot n_{\cdot jk}}\right)\]
Di bawah hipotesis nol independensi kondisional, \(G^2\) mengikuti
distribusi \(\chi^2\) dengan derajat bebas \((I-1)(J-1)|K|\) di mana
\(|K|\) adalah jumlah kombinasi nilai variabel kondisional.

\begin{Shaded}
\begin{Highlighting}[]
\CommentTok{\# Salin data imputed untuk proses diskretisasi}
\NormalTok{diabetes\_disc\_prep }\OtherTok{\textless{}{-}}\NormalTok{ diabetes\_imp}

\CommentTok{\# Konversi Outcome ke faktor terlebih dahulu (agar tidak ikut didiskretisasi)}
\NormalTok{diabetes\_disc\_prep}\SpecialCharTok{$}\NormalTok{Outcome }\OtherTok{\textless{}{-}} \FunctionTok{factor}\NormalTok{(}
\NormalTok{  diabetes\_disc\_prep}\SpecialCharTok{$}\NormalTok{Outcome,}
  \AttributeTok{levels =} \FunctionTok{c}\NormalTok{(}\DecValTok{0}\NormalTok{, }\DecValTok{1}\NormalTok{),}
  \AttributeTok{labels =} \FunctionTok{c}\NormalTok{(}\StringTok{"Negatif"}\NormalTok{, }\StringTok{"Positif"}\NormalTok{)}
\NormalTok{)}

\CommentTok{\# Diskretisasi semua variabel numerik menjadi 3 bin (kuantil)}
\CommentTok{\# Menggunakan metode \textquotesingle{}interval\textquotesingle{} untuk menghindari error zero{-}length intervals}
\CommentTok{\# akibat banyaknya data yang bernilai persis sama hasil imputasi median.}
\NormalTok{disc\_data }\OtherTok{\textless{}{-}}\NormalTok{ bnlearn}\SpecialCharTok{::}\FunctionTok{discretize}\NormalTok{(diabetes\_disc\_prep,}
                                  \AttributeTok{method =} \StringTok{"interval"}\NormalTok{,}
                                  \AttributeTok{breaks  =} \DecValTok{3}\NormalTok{)}

\CommentTok{\# Tampilkan tipe data setelah diskretisasi dalam format tabel}
\NormalTok{tipe\_disc\_df }\OtherTok{\textless{}{-}} \FunctionTok{data.frame}\NormalTok{(}
  \AttributeTok{Variabel =} \FunctionTok{names}\NormalTok{(disc\_data),}
  \AttributeTok{Tipe\_Data =} \FunctionTok{sapply}\NormalTok{(disc\_data, class),}
  \AttributeTok{row.names =} \ConstantTok{NULL}
\NormalTok{)}
\NormalTok{knitr}\SpecialCharTok{::}\FunctionTok{kable}\NormalTok{(tipe\_disc\_df, }\AttributeTok{caption =} \StringTok{"Tipe Data Setelah Diskretisasi"}\NormalTok{)}
\end{Highlighting}
\end{Shaded}

\begin{longtable}[]{@{}ll@{}}
\caption{Tipe Data Setelah Diskretisasi}\tabularnewline
\toprule\noalign{}
Variabel & Tipe\_Data \\
\midrule\noalign{}
\endfirsthead
\toprule\noalign{}
Variabel & Tipe\_Data \\
\midrule\noalign{}
\endhead
\bottomrule\noalign{}
\endlastfoot
Pregnancies & factor \\
Glucose & factor \\
BloodPressure & factor \\
SkinThickness & factor \\
Insulin & factor \\
BMI & factor \\
DiabetesPedigreeFunction & factor \\
Age & factor \\
Outcome & factor \\
\end{longtable}

\begin{Shaded}
\begin{Highlighting}[]
\CommentTok{\# Tampilkan level dalam bentuk tabel agar rapi (tidak berantakan di layar)}
\NormalTok{levels\_df }\OtherTok{\textless{}{-}} \FunctionTok{data.frame}\NormalTok{(}
  \AttributeTok{Variabel =} \FunctionTok{names}\NormalTok{(disc\_data),}
  \AttributeTok{Level\_Kategori =} \FunctionTok{sapply}\NormalTok{(}\FunctionTok{names}\NormalTok{(disc\_data), }\ControlFlowTok{function}\NormalTok{(col) }\FunctionTok{paste}\NormalTok{(}\FunctionTok{levels}\NormalTok{(disc\_data[[col]]), }\AttributeTok{collapse =} \StringTok{" | "}\NormalTok{)),}
  \AttributeTok{row.names =} \ConstantTok{NULL}
\NormalTok{)}
\NormalTok{knitr}\SpecialCharTok{::}\FunctionTok{kable}\NormalTok{(levels\_df, }\AttributeTok{caption =} \StringTok{"Level Variabel Setelah Diskretisasi"}\NormalTok{)}
\end{Highlighting}
\end{Shaded}

\begin{longtable}[]{@{}
  >{\raggedright\arraybackslash}p{(\linewidth - 2\tabcolsep) * \real{0.2778}}
  >{\raggedright\arraybackslash}p{(\linewidth - 2\tabcolsep) * \real{0.7222}}@{}}
\caption{Level Variabel Setelah Diskretisasi}\tabularnewline
\toprule\noalign{}
\begin{minipage}[b]{\linewidth}\raggedright
Variabel
\end{minipage} & \begin{minipage}[b]{\linewidth}\raggedright
Level\_Kategori
\end{minipage} \\
\midrule\noalign{}
\endfirsthead
\toprule\noalign{}
\begin{minipage}[b]{\linewidth}\raggedright
Variabel
\end{minipage} & \begin{minipage}[b]{\linewidth}\raggedright
Level\_Kategori
\end{minipage} \\
\midrule\noalign{}
\endhead
\bottomrule\noalign{}
\endlastfoot
Pregnancies & {[}0,5.66667{]} \textbar{} (5.66667,11.3333{]} \textbar{}
(11.3333,17{]} \\
Glucose & {[}44,95.6667{]} \textbar{} (95.6667,147.333{]} \textbar{}
(147.333,199{]} \\
BloodPressure & {[}24,56.6667{]} \textbar{} (56.6667,89.3333{]}
\textbar{} (89.3333,122{]} \\
SkinThickness & {[}7,37.6667{]} \textbar{} (37.6667,68.3333{]}
\textbar{} (68.3333,99{]} \\
Insulin & {[}14,291.333{]} \textbar{} (291.333,568.667{]} \textbar{}
(568.667,846{]} \\
BMI & {[}18.2,34.5{]} \textbar{} (34.5,50.8{]} \textbar{}
(50.8,67.1{]} \\
DiabetesPedigreeFunction & {[}0.078,0.858667{]} \textbar{}
(0.858667,1.63933{]} \textbar{} (1.63933,2.42{]} \\
Age & {[}21,41{]} \textbar{} (41,61{]} \textbar{} (61,81{]} \\
Outcome & Negatif \textbar{} Positif \\
\end{longtable}

\subsection{Algoritma PC-Stable}\label{algoritma-pc-stable}

Algoritma PC-Stable (Colombo \& Maathuis, 2014) adalah varian
\emph{order-independent} dari PC Algorithm yang bekerja dalam dua fase
utama:

\begin{enumerate}
\def\labelenumi{\arabic{enumi}.}
\item
  \textbf{Fase Skeleton:} Dimulai dari graf lengkap (\emph{complete
  graph}), setiap pasang variabel \((X, Y)\) diuji apakah saling bebas
  secara kondisional diberikan berbagai himpunan \(S\) (dimulai dari
  \(|S|=0, 1, 2, \ldots\)). Jika ditemukan \(S\) (\emph{separating set})
  sehingga \(X \perp\!\!\!\perp Y \mid S\), maka edge \(X-Y\) dihapus.
\item
  \textbf{Fase Orientasi:} Edge yang tersisa diorientasikan menggunakan
  \emph{Meek Rules}: orientasi dimulai dari deteksi \emph{v-structures}
  (unshielded colliders \(X \to Z \leftarrow Y\) di mana \(X-Y\) tidak
  terhubung langsung), dilanjutkan dengan penerapan empat orientation
  rules untuk menghindari siklus dan v-structure baru.
\end{enumerate}

\textbf{Parameter yang digunakan:} - CI Test: \texttt{"mi"} (G-Squared /
Mutual Information) - sesuai untuk data diskrit - Nilai \(\alpha\): 0.01
dan 0.05 (\emph{sensitivity analysis}) - Paket: \texttt{bnlearn} (versi
\(\geq\) 4.7) - Platform: R \(\geq\) 4.1.0

\begin{center}\rule{0.5\linewidth}{0.5pt}\end{center}

\section{Hasil}\label{hasil}

\subsection{Statistik Deskriptif Pasca
Preprocessing}\label{statistik-deskriptif-pasca-preprocessing}

\begin{Shaded}
\begin{Highlighting}[]
\CommentTok{\# Transpose tabel agar rapi dan muat di halaman}
\NormalTok{post\_stats }\OtherTok{\textless{}{-}} \FunctionTok{t}\NormalTok{(}\FunctionTok{round}\NormalTok{(}\FunctionTok{sapply}\NormalTok{(diabetes\_imp[, }\SpecialCharTok{{-}}\DecValTok{9}\NormalTok{], }\ControlFlowTok{function}\NormalTok{(x) \{}
  \FunctionTok{c}\NormalTok{(}\AttributeTok{Min    =} \FunctionTok{min}\NormalTok{(x),}
    \AttributeTok{Q1     =} \FunctionTok{quantile}\NormalTok{(x, }\FloatTok{0.25}\NormalTok{),}
    \AttributeTok{Median =} \FunctionTok{median}\NormalTok{(x),}
    \AttributeTok{Mean   =} \FunctionTok{mean}\NormalTok{(x),}
    \AttributeTok{Q3     =} \FunctionTok{quantile}\NormalTok{(x, }\FloatTok{0.75}\NormalTok{),}
    \AttributeTok{Max    =} \FunctionTok{max}\NormalTok{(x),}
    \AttributeTok{SD     =} \FunctionTok{sd}\NormalTok{(x))}
\NormalTok{\}), }\DecValTok{2}\NormalTok{))}
\NormalTok{knitr}\SpecialCharTok{::}\FunctionTok{kable}\NormalTok{(post\_stats, }\AttributeTok{caption =} \StringTok{"Statistik Deskriptif Setelah Imputasi (Variabel Kontinu)"}\NormalTok{, }\AttributeTok{row.names =} \ConstantTok{TRUE}\NormalTok{)}
\end{Highlighting}
\end{Shaded}

\begin{longtable}[]{@{}
  >{\raggedright\arraybackslash}p{(\linewidth - 14\tabcolsep) * \real{0.3472}}
  >{\raggedleft\arraybackslash}p{(\linewidth - 14\tabcolsep) * \real{0.0833}}
  >{\raggedleft\arraybackslash}p{(\linewidth - 14\tabcolsep) * \real{0.0972}}
  >{\raggedleft\arraybackslash}p{(\linewidth - 14\tabcolsep) * \real{0.0972}}
  >{\raggedleft\arraybackslash}p{(\linewidth - 14\tabcolsep) * \real{0.0972}}
  >{\raggedleft\arraybackslash}p{(\linewidth - 14\tabcolsep) * \real{0.0972}}
  >{\raggedleft\arraybackslash}p{(\linewidth - 14\tabcolsep) * \real{0.0972}}
  >{\raggedleft\arraybackslash}p{(\linewidth - 14\tabcolsep) * \real{0.0833}}@{}}
\caption{Statistik Deskriptif Setelah Imputasi (Variabel
Kontinu)}\tabularnewline
\toprule\noalign{}
\begin{minipage}[b]{\linewidth}\raggedright
\end{minipage} & \begin{minipage}[b]{\linewidth}\raggedleft
Min
\end{minipage} & \begin{minipage}[b]{\linewidth}\raggedleft
Q1.25\%
\end{minipage} & \begin{minipage}[b]{\linewidth}\raggedleft
Median
\end{minipage} & \begin{minipage}[b]{\linewidth}\raggedleft
Mean
\end{minipage} & \begin{minipage}[b]{\linewidth}\raggedleft
Q3.75\%
\end{minipage} & \begin{minipage}[b]{\linewidth}\raggedleft
Max
\end{minipage} & \begin{minipage}[b]{\linewidth}\raggedleft
SD
\end{minipage} \\
\midrule\noalign{}
\endfirsthead
\toprule\noalign{}
\begin{minipage}[b]{\linewidth}\raggedright
\end{minipage} & \begin{minipage}[b]{\linewidth}\raggedleft
Min
\end{minipage} & \begin{minipage}[b]{\linewidth}\raggedleft
Q1.25\%
\end{minipage} & \begin{minipage}[b]{\linewidth}\raggedleft
Median
\end{minipage} & \begin{minipage}[b]{\linewidth}\raggedleft
Mean
\end{minipage} & \begin{minipage}[b]{\linewidth}\raggedleft
Q3.75\%
\end{minipage} & \begin{minipage}[b]{\linewidth}\raggedleft
Max
\end{minipage} & \begin{minipage}[b]{\linewidth}\raggedleft
SD
\end{minipage} \\
\midrule\noalign{}
\endhead
\bottomrule\noalign{}
\endlastfoot
Pregnancies & 0.00 & 1.00 & 3.00 & 3.85 & 6.00 & 17.00 & 3.37 \\
Glucose & 44.00 & 99.75 & 117.00 & 121.66 & 140.25 & 199.00 & 30.44 \\
BloodPressure & 24.00 & 64.00 & 72.00 & 72.39 & 80.00 & 122.00 &
12.10 \\
SkinThickness & 7.00 & 25.00 & 29.00 & 29.11 & 32.00 & 99.00 & 8.79 \\
Insulin & 14.00 & 121.50 & 125.00 & 140.67 & 127.25 & 846.00 & 86.38 \\
BMI & 18.20 & 27.50 & 32.30 & 32.46 & 36.60 & 67.10 & 6.88 \\
DiabetesPedigreeFunction & 0.08 & 0.24 & 0.37 & 0.47 & 0.63 & 2.42 &
0.33 \\
Age & 21.00 & 24.00 & 29.00 & 33.24 & 41.00 & 81.00 & 11.76 \\
\end{longtable}

\begin{Shaded}
\begin{Highlighting}[]
\NormalTok{outcome\_counts }\OtherTok{\textless{}{-}} \FunctionTok{table}\NormalTok{(diabetes\_imp}\SpecialCharTok{$}\NormalTok{Outcome)}
\NormalTok{outcome\_pct    }\OtherTok{\textless{}{-}} \FunctionTok{round}\NormalTok{(}\FunctionTok{prop.table}\NormalTok{(outcome\_counts) }\SpecialCharTok{*} \DecValTok{100}\NormalTok{, }\DecValTok{1}\NormalTok{)}
\FunctionTok{barplot}\NormalTok{(outcome\_counts,}
        \AttributeTok{col     =} \FunctionTok{c}\NormalTok{(}\StringTok{"\#AED6F1"}\NormalTok{, }\StringTok{"\#E74C3C"}\NormalTok{),}
        \AttributeTok{names.arg =} \FunctionTok{c}\NormalTok{(}\FunctionTok{paste0}\NormalTok{(}\StringTok{"Negatif (0)}\SpecialCharTok{\textbackslash{}n}\StringTok{"}\NormalTok{, outcome\_pct[}\DecValTok{1}\NormalTok{], }\StringTok{"\%"}\NormalTok{),}
                      \FunctionTok{paste0}\NormalTok{(}\StringTok{"Positif (1)}\SpecialCharTok{\textbackslash{}n}\StringTok{"}\NormalTok{,  outcome\_pct[}\DecValTok{2}\NormalTok{], }\StringTok{"\%"}\NormalTok{)),}
        \AttributeTok{main    =} \StringTok{"Distribusi Outcome (Diagnosis Diabetes)"}\NormalTok{,}
        \AttributeTok{ylab    =} \StringTok{"Frekuensi"}\NormalTok{,}
        \AttributeTok{border  =} \StringTok{"white"}\NormalTok{)}
\FunctionTok{legend}\NormalTok{(}\StringTok{"topright"}\NormalTok{,}
       \AttributeTok{legend =} \FunctionTok{c}\NormalTok{(}\FunctionTok{paste0}\NormalTok{(}\StringTok{"Negatif: n = "}\NormalTok{, outcome\_counts[}\DecValTok{1}\NormalTok{]),}
                  \FunctionTok{paste0}\NormalTok{(}\StringTok{"Positif: n = "}\NormalTok{, outcome\_counts[}\DecValTok{2}\NormalTok{])),}
       \AttributeTok{fill   =} \FunctionTok{c}\NormalTok{(}\StringTok{"\#AED6F1"}\NormalTok{, }\StringTok{"\#E74C3C"}\NormalTok{), }\AttributeTok{bty =} \StringTok{"n"}\NormalTok{)}
\end{Highlighting}
\end{Shaded}

\begin{figure}

{\centering \includegraphics[width=0.85\linewidth]{pc_algorithm_vs_ges_files/figure-latex/outcome-bar-1} 

}

\caption{Distribusi Variabel Outcome (Diagnosis Diabetes)}\label{fig:outcome-bar}
\end{figure}

\begin{Shaded}
\begin{Highlighting}[]
\NormalTok{cor\_mat }\OtherTok{\textless{}{-}} \FunctionTok{cor}\NormalTok{(diabetes\_imp[, }\SpecialCharTok{{-}}\DecValTok{9}\NormalTok{])}
\FunctionTok{corrplot}\NormalTok{(cor\_mat,}
         \AttributeTok{method      =} \StringTok{"color"}\NormalTok{,}
         \AttributeTok{type        =} \StringTok{"upper"}\NormalTok{,}
         \AttributeTok{addCoef.col =} \StringTok{"black"}\NormalTok{,}
         \AttributeTok{number.cex  =} \FloatTok{0.65}\NormalTok{,}
         \AttributeTok{tl.cex      =} \FloatTok{0.75}\NormalTok{,}
         \AttributeTok{col         =} \FunctionTok{colorRampPalette}\NormalTok{(}\FunctionTok{c}\NormalTok{(}\StringTok{"\#D73027"}\NormalTok{, }\StringTok{"\#FFFFFF"}\NormalTok{, }\StringTok{"\#1A9850"}\NormalTok{))(}\DecValTok{200}\NormalTok{),}
         \AttributeTok{title       =} \StringTok{"Matriks Korelasi (Variabel Kontinu)"}\NormalTok{,}
         \AttributeTok{mar         =} \FunctionTok{c}\NormalTok{(}\DecValTok{0}\NormalTok{, }\DecValTok{0}\NormalTok{, }\DecValTok{2}\NormalTok{, }\DecValTok{0}\NormalTok{))}
\end{Highlighting}
\end{Shaded}

\begin{figure}

{\centering \includegraphics[width=0.85\linewidth]{pc_algorithm_vs_ges_files/figure-latex/corrplot-1} 

}

\caption{Matriks Korelasi Pearson Antar Variabel Numerik (Pasca Imputasi)}\label{fig:corrplot}
\end{figure}

\subsection{Causal Discovery: Algoritma
PC-Stable}\label{causal-discovery-algoritma-pc-stable}

\subsubsection{Fungsi Utilitas Visualisasi
CPDAG}\label{fungsi-utilitas-visualisasi-cpdag}

\begin{Shaded}
\begin{Highlighting}[]
\CommentTok{\# Fungsi plot CPDAG menggunakan igraph}
\CommentTok{\# Edge berarah  {-}\textgreater{} ditampilkan dengan warna gelap dan anak panah}
\CommentTok{\# Edge tidak berarah {-} ditampilkan merah putus{-}putus tanpa anak panah}
\NormalTok{plot\_cpdag }\OtherTok{\textless{}{-}} \ControlFlowTok{function}\NormalTok{(net, }\AttributeTok{title =} \StringTok{"CPDAG"}\NormalTok{) \{}

\NormalTok{  node\_names }\OtherTok{\textless{}{-}} \FunctionTok{nodes}\NormalTok{(net)}
\NormalTok{  n          }\OtherTok{\textless{}{-}} \FunctionTok{length}\NormalTok{(node\_names)}
\NormalTok{  dir\_a      }\OtherTok{\textless{}{-}} \FunctionTok{directed.arcs}\NormalTok{(net)}
\NormalTok{  undir\_a    }\OtherTok{\textless{}{-}} \FunctionTok{undirected.arcs}\NormalTok{(net)}

  \CommentTok{\# Akumulasi daftar edge}
\NormalTok{  e\_from }\OtherTok{\textless{}{-}} \FunctionTok{integer}\NormalTok{(}\DecValTok{0}\NormalTok{); e\_to }\OtherTok{\textless{}{-}} \FunctionTok{integer}\NormalTok{(}\DecValTok{0}\NormalTok{)}
\NormalTok{  e\_col  }\OtherTok{\textless{}{-}} \FunctionTok{character}\NormalTok{(}\DecValTok{0}\NormalTok{); e\_lty }\OtherTok{\textless{}{-}} \FunctionTok{integer}\NormalTok{(}\DecValTok{0}\NormalTok{); e\_arr }\OtherTok{\textless{}{-}} \FunctionTok{integer}\NormalTok{(}\DecValTok{0}\NormalTok{)}

  \CommentTok{\# Edge berarah}
  \ControlFlowTok{if}\NormalTok{ (}\FunctionTok{nrow}\NormalTok{(dir\_a) }\SpecialCharTok{\textgreater{}} \DecValTok{0}\NormalTok{) \{}
\NormalTok{    e\_from }\OtherTok{\textless{}{-}} \FunctionTok{c}\NormalTok{(e\_from, }\FunctionTok{match}\NormalTok{(dir\_a[, }\StringTok{"from"}\NormalTok{], node\_names))}
\NormalTok{    e\_to   }\OtherTok{\textless{}{-}} \FunctionTok{c}\NormalTok{(e\_to,   }\FunctionTok{match}\NormalTok{(dir\_a[, }\StringTok{"to"}\NormalTok{],   node\_names))}
\NormalTok{    e\_col  }\OtherTok{\textless{}{-}} \FunctionTok{c}\NormalTok{(e\_col,  }\FunctionTok{rep}\NormalTok{(}\StringTok{"\#1A5276"}\NormalTok{, }\FunctionTok{nrow}\NormalTok{(dir\_a)))}
\NormalTok{    e\_lty  }\OtherTok{\textless{}{-}} \FunctionTok{c}\NormalTok{(e\_lty,  }\FunctionTok{rep}\NormalTok{(}\DecValTok{1}\DataTypeTok{L}\NormalTok{,        }\FunctionTok{nrow}\NormalTok{(dir\_a)))}
\NormalTok{    e\_arr  }\OtherTok{\textless{}{-}} \FunctionTok{c}\NormalTok{(e\_arr,  }\FunctionTok{rep}\NormalTok{(}\DecValTok{2}\DataTypeTok{L}\NormalTok{,        }\FunctionTok{nrow}\NormalTok{(dir\_a)))   }\CommentTok{\# anak panah ke depan}
\NormalTok{  \}}

  \CommentTok{\# Edge tidak berarah {-} ambil pasangan unik saja}
  \ControlFlowTok{if}\NormalTok{ (}\FunctionTok{nrow}\NormalTok{(undir\_a) }\SpecialCharTok{\textgreater{}} \DecValTok{0}\NormalTok{) \{}
\NormalTok{    pairs\_m }\OtherTok{\textless{}{-}} \FunctionTok{unique}\NormalTok{(}\FunctionTok{t}\NormalTok{(}\FunctionTok{apply}\NormalTok{(undir\_a, }\DecValTok{1}\DataTypeTok{L}\NormalTok{, sort)))}
    \ControlFlowTok{if}\NormalTok{ (}\SpecialCharTok{!}\FunctionTok{is.matrix}\NormalTok{(pairs\_m)) pairs\_m }\OtherTok{\textless{}{-}} \FunctionTok{matrix}\NormalTok{(pairs\_m, }\AttributeTok{nrow =} \DecValTok{1}\DataTypeTok{L}\NormalTok{)}
\NormalTok{    e\_from }\OtherTok{\textless{}{-}} \FunctionTok{c}\NormalTok{(e\_from, }\FunctionTok{match}\NormalTok{(pairs\_m[, }\DecValTok{1}\DataTypeTok{L}\NormalTok{], node\_names))}
\NormalTok{    e\_to   }\OtherTok{\textless{}{-}} \FunctionTok{c}\NormalTok{(e\_to,   }\FunctionTok{match}\NormalTok{(pairs\_m[, }\DecValTok{2}\DataTypeTok{L}\NormalTok{], node\_names))}
\NormalTok{    e\_col  }\OtherTok{\textless{}{-}} \FunctionTok{c}\NormalTok{(e\_col,  }\FunctionTok{rep}\NormalTok{(}\StringTok{"\#C0392B"}\NormalTok{, }\FunctionTok{nrow}\NormalTok{(pairs\_m)))}
\NormalTok{    e\_lty  }\OtherTok{\textless{}{-}} \FunctionTok{c}\NormalTok{(e\_lty,  }\FunctionTok{rep}\NormalTok{(}\DecValTok{2}\DataTypeTok{L}\NormalTok{,        }\FunctionTok{nrow}\NormalTok{(pairs\_m)))}
\NormalTok{    e\_arr  }\OtherTok{\textless{}{-}} \FunctionTok{c}\NormalTok{(e\_arr,  }\FunctionTok{rep}\NormalTok{(}\DecValTok{0}\DataTypeTok{L}\NormalTok{,        }\FunctionTok{nrow}\NormalTok{(pairs\_m)))   }\CommentTok{\# tanpa anak panah}
\NormalTok{  \}}

  \CommentTok{\# Buat igraph}
\NormalTok{  g }\OtherTok{\textless{}{-}} \FunctionTok{make\_empty\_graph}\NormalTok{(}\AttributeTok{n =}\NormalTok{ n, }\AttributeTok{directed =} \ConstantTok{TRUE}\NormalTok{)}
  \FunctionTok{V}\NormalTok{(g)}\SpecialCharTok{$}\NormalTok{name }\OtherTok{\textless{}{-}}\NormalTok{ node\_names}

  \ControlFlowTok{if}\NormalTok{ (}\FunctionTok{length}\NormalTok{(e\_from) }\SpecialCharTok{\textgreater{}} \DecValTok{0}\NormalTok{) \{}
\NormalTok{    g               }\OtherTok{\textless{}{-}} \FunctionTok{add\_edges}\NormalTok{(g, }\FunctionTok{as.vector}\NormalTok{(}\FunctionTok{rbind}\NormalTok{(e\_from, e\_to)))}
    \FunctionTok{E}\NormalTok{(g)}\SpecialCharTok{$}\NormalTok{color      }\OtherTok{\textless{}{-}}\NormalTok{ e\_col}
    \FunctionTok{E}\NormalTok{(g)}\SpecialCharTok{$}\NormalTok{lty        }\OtherTok{\textless{}{-}}\NormalTok{ e\_lty}
    \FunctionTok{E}\NormalTok{(g)}\SpecialCharTok{$}\NormalTok{arrow.mode }\OtherTok{\textless{}{-}}\NormalTok{ e\_arr}
\NormalTok{  \}}

  \FunctionTok{set.seed}\NormalTok{(}\DecValTok{123}\NormalTok{)}
\NormalTok{  lo }\OtherTok{\textless{}{-}} \FunctionTok{layout\_with\_fr}\NormalTok{(g)}

  \FunctionTok{plot}\NormalTok{(g,}
       \AttributeTok{layout             =}\NormalTok{ lo,}
       \AttributeTok{vertex.color       =} \StringTok{"\#AED6F1"}\NormalTok{,}
       \AttributeTok{vertex.frame.color =} \StringTok{"\#1A5276"}\NormalTok{,}
       \AttributeTok{vertex.shape       =} \StringTok{"circle"}\NormalTok{,}
       \AttributeTok{vertex.size        =} \DecValTok{30}\NormalTok{,}
       \AttributeTok{vertex.label       =} \FunctionTok{V}\NormalTok{(g)}\SpecialCharTok{$}\NormalTok{name,}
       \AttributeTok{vertex.label.cex   =} \FloatTok{0.72}\NormalTok{,}
       \AttributeTok{vertex.label.color =} \StringTok{"black"}\NormalTok{,}
       \AttributeTok{vertex.label.font  =} \DecValTok{2}\DataTypeTok{L}\NormalTok{,}
       \AttributeTok{edge.arrow.size    =} \FloatTok{0.55}\NormalTok{,}
       \AttributeTok{edge.width         =} \FloatTok{1.8}\NormalTok{,}
       \AttributeTok{main               =}\NormalTok{ title)}

  \FunctionTok{legend}\NormalTok{(}\StringTok{"bottomright"}\NormalTok{,}
         \AttributeTok{legend =} \FunctionTok{c}\NormalTok{(}\StringTok{"Edge Berarah ({-}\textgreater{})"}\NormalTok{, }\StringTok{"Edge Tidak Berarah ({-})"}\NormalTok{),}
         \AttributeTok{col    =} \FunctionTok{c}\NormalTok{(}\StringTok{"\#1A5276"}\NormalTok{, }\StringTok{"\#C0392B"}\NormalTok{),}
         \AttributeTok{lty    =} \FunctionTok{c}\NormalTok{(}\DecValTok{1}\DataTypeTok{L}\NormalTok{, }\DecValTok{2}\DataTypeTok{L}\NormalTok{),}
         \AttributeTok{lwd    =} \DecValTok{2}\NormalTok{,}
         \AttributeTok{cex    =} \FloatTok{0.75}\NormalTok{,}
         \AttributeTok{bg     =} \StringTok{"white"}\NormalTok{,}
         \AttributeTok{box.lty =} \DecValTok{0}\NormalTok{)}
\NormalTok{\}}

\CommentTok{\# Fungsi ringkasan edge}
\NormalTok{summarize\_edges }\OtherTok{\textless{}{-}} \ControlFlowTok{function}\NormalTok{(net, }\AttributeTok{label =} \StringTok{""}\NormalTok{) \{}
\NormalTok{  dir\_a   }\OtherTok{\textless{}{-}} \FunctionTok{directed.arcs}\NormalTok{(net)}
\NormalTok{  undir\_a }\OtherTok{\textless{}{-}} \FunctionTok{undirected.arcs}\NormalTok{(net)}
\NormalTok{  n\_dir   }\OtherTok{\textless{}{-}} \FunctionTok{nrow}\NormalTok{(dir\_a)}
\NormalTok{  n\_undir }\OtherTok{\textless{}{-}} \ControlFlowTok{if}\NormalTok{ (}\FunctionTok{nrow}\NormalTok{(undir\_a) }\SpecialCharTok{\textgreater{}} \DecValTok{0}\NormalTok{) }\FunctionTok{nrow}\NormalTok{(undir\_a) }\SpecialCharTok{/} \DecValTok{2}\DataTypeTok{L} \ControlFlowTok{else} \DecValTok{0}\DataTypeTok{L}
  \FunctionTok{cat}\NormalTok{(}\FunctionTok{sprintf}\NormalTok{(}\StringTok{"}\SpecialCharTok{\textbackslash{}n}\StringTok{=== \%s ===}\SpecialCharTok{\textbackslash{}n}\StringTok{"}\NormalTok{, label))}
  \FunctionTok{cat}\NormalTok{(}\FunctionTok{sprintf}\NormalTok{(}\StringTok{"Edge berarah    : \%d}\SpecialCharTok{\textbackslash{}n}\StringTok{"}\NormalTok{, n\_dir))}
  \FunctionTok{cat}\NormalTok{(}\FunctionTok{sprintf}\NormalTok{(}\StringTok{"Edge tdk berarah: \%d}\SpecialCharTok{\textbackslash{}n}\StringTok{"}\NormalTok{, n\_undir))}
  \FunctionTok{cat}\NormalTok{(}\FunctionTok{sprintf}\NormalTok{(}\StringTok{"Total edge      : \%d}\SpecialCharTok{\textbackslash{}n\textbackslash{}n}\StringTok{"}\NormalTok{, n\_dir }\SpecialCharTok{+}\NormalTok{ n\_undir))}

  \ControlFlowTok{if}\NormalTok{ (n\_dir }\SpecialCharTok{\textgreater{}} \DecValTok{0}\NormalTok{) \{}
    \FunctionTok{cat}\NormalTok{(}\StringTok{"Edge berarah:}\SpecialCharTok{\textbackslash{}n}\StringTok{"}\NormalTok{); }\FunctionTok{print}\NormalTok{(dir\_a)}
\NormalTok{  \}}
  \ControlFlowTok{if}\NormalTok{ (n\_undir }\SpecialCharTok{\textgreater{}} \DecValTok{0}\NormalTok{) \{}
\NormalTok{    pairs\_m }\OtherTok{\textless{}{-}} \FunctionTok{unique}\NormalTok{(}\FunctionTok{t}\NormalTok{(}\FunctionTok{apply}\NormalTok{(undir\_a, }\DecValTok{1}\DataTypeTok{L}\NormalTok{, sort)))}
    \ControlFlowTok{if}\NormalTok{ (}\SpecialCharTok{!}\FunctionTok{is.matrix}\NormalTok{(pairs\_m)) pairs\_m }\OtherTok{\textless{}{-}} \FunctionTok{matrix}\NormalTok{(pairs\_m, }\AttributeTok{nrow =} \DecValTok{1}\DataTypeTok{L}\NormalTok{)}
    \FunctionTok{cat}\NormalTok{(}\StringTok{"}\SpecialCharTok{\textbackslash{}n}\StringTok{Edge tidak berarah (pasangan unik):}\SpecialCharTok{\textbackslash{}n}\StringTok{"}\NormalTok{); }\FunctionTok{print}\NormalTok{(pairs\_m)}
\NormalTok{  \}}
  \FunctionTok{invisible}\NormalTok{(}\FunctionTok{list}\NormalTok{(}\AttributeTok{n\_dir =}\NormalTok{ n\_dir, }\AttributeTok{n\_undir =}\NormalTok{ n\_undir))}
\NormalTok{\}}
\end{Highlighting}
\end{Shaded}

\subsubsection{\texorpdfstring{PC-Stable dengan
\(\alpha = 0.01\)}{PC-Stable dengan \textbackslash alpha = 0.01}}\label{pc-stable-dengan-alpha-0.01}

\begin{Shaded}
\begin{Highlighting}[]
\FunctionTok{set.seed}\NormalTok{(}\DecValTok{42}\NormalTok{)}
\NormalTok{pc\_01  }\OtherTok{\textless{}{-}} \FunctionTok{pc.stable}\NormalTok{(disc\_data, }\AttributeTok{alpha =} \FloatTok{0.01}\NormalTok{, }\AttributeTok{test =} \StringTok{"mi"}\NormalTok{, }\AttributeTok{undirected =} \ConstantTok{FALSE}\NormalTok{)}
\NormalTok{info01 }\OtherTok{\textless{}{-}} \FunctionTok{summarize\_edges}\NormalTok{(pc\_01, }\StringTok{"PC{-}Stable (alpha = 0.01)"}\NormalTok{)}
\end{Highlighting}
\end{Shaded}

\begin{verbatim}
## 
## === PC-Stable (alpha = 0.01) ===
## Edge berarah    : 2
## Edge tdk berarah: 2
## Total edge      : 4
## 
## Edge berarah:
##      from      to       
## [1,] "Insulin" "Glucose"
## [2,] "Outcome" "Glucose"
## 
## Edge tidak berarah (pasangan unik):
##      [,1]  [,2]           
## [1,] "Age" "Pregnancies"  
## [2,] "BMI" "SkinThickness"
\end{verbatim}

\begin{Shaded}
\begin{Highlighting}[]
\FunctionTok{plot\_cpdag}\NormalTok{(pc\_01, }\AttributeTok{title =} \FunctionTok{expression}\NormalTok{(}\FunctionTok{paste}\NormalTok{(}\StringTok{"CPDAG: PC{-}Stable  ("}\NormalTok{, alpha, }\StringTok{" = 0.01)"}\NormalTok{)))}
\end{Highlighting}
\end{Shaded}

\begin{figure}

{\centering \includegraphics[width=0.85\linewidth]{pc_algorithm_vs_ges_files/figure-latex/pc-alpha01-1} 

}

\caption{CPDAG Hasil PC-Stable (alpha = 0.01)}\label{fig:pc-alpha01}
\end{figure}

\subsubsection{\texorpdfstring{PC-Stable dengan
\(\alpha = 0.05\)}{PC-Stable dengan \textbackslash alpha = 0.05}}\label{pc-stable-dengan-alpha-0.05}

\begin{Shaded}
\begin{Highlighting}[]
\FunctionTok{set.seed}\NormalTok{(}\DecValTok{42}\NormalTok{)}
\NormalTok{pc\_05  }\OtherTok{\textless{}{-}} \FunctionTok{pc.stable}\NormalTok{(disc\_data, }\AttributeTok{alpha =} \FloatTok{0.05}\NormalTok{, }\AttributeTok{test =} \StringTok{"mi"}\NormalTok{, }\AttributeTok{undirected =} \ConstantTok{FALSE}\NormalTok{)}
\NormalTok{info05 }\OtherTok{\textless{}{-}} \FunctionTok{summarize\_edges}\NormalTok{(pc\_05, }\StringTok{"PC{-}Stable (alpha = 0.05)"}\NormalTok{)}
\end{Highlighting}
\end{Shaded}

\begin{verbatim}
## 
## === PC-Stable (alpha = 0.05) ===
## Edge berarah    : 2
## Edge tdk berarah: 3
## Total edge      : 5
## 
## Edge berarah:
##      from            to   
## [1,] "Pregnancies"   "Age"
## [2,] "BloodPressure" "Age"
## 
## Edge tidak berarah (pasangan unik):
##      [,1]      [,2]           
## [1,] "Glucose" "Insulin"      
## [2,] "Glucose" "Outcome"      
## [3,] "BMI"     "SkinThickness"
\end{verbatim}

\begin{Shaded}
\begin{Highlighting}[]
\FunctionTok{plot\_cpdag}\NormalTok{(pc\_05, }\AttributeTok{title =} \FunctionTok{expression}\NormalTok{(}\FunctionTok{paste}\NormalTok{(}\StringTok{"CPDAG: PC{-}Stable  ("}\NormalTok{, alpha, }\StringTok{" = 0.05)"}\NormalTok{)))}
\end{Highlighting}
\end{Shaded}

\begin{figure}

{\centering \includegraphics[width=0.85\linewidth]{pc_algorithm_vs_ges_files/figure-latex/pc-alpha05-1} 

}

\caption{CPDAG Hasil PC-Stable (alpha = 0.05)}\label{fig:pc-alpha05}
\end{figure}

\subsubsection{Score-Based Algorithm: Hill-Climbing (Alternatif
GES)}\label{score-based-algorithm-hill-climbing-alternatif-ges}

Algoritma \emph{Score-Based} bekerja dengan mengeksplorasi ruang
pencarian DAG secara rakus (\emph{greedy search}) untuk menemukan graf
yang memaksimalkan fungsi skor (\emph{goodness-of-fit}). Dalam paket
\texttt{bnlearn}, algoritma utama untuk pendekatan ini adalah
\textbf{Hill-Climbing (\texttt{hc})}, yang filosofinya sangat mirip
dengan \emph{Greedy Equivalence Search} (GES).

\textbf{Parameter yang digunakan:} - \textbf{Skor:} \texttt{"bic"}
(Bayesian Information Criterion). Pemilihan BIC dijustifikasi karena
memberikan penalti pada kompleksitas graf (mencegah \emph{overfitting}
yang sering terjadi jika hanya menggunakan \emph{log-likelihood}). -
Graf hasil \emph{Hill-Climbing} berupa DAG kemudian dikonversi menjadi
\emph{CPDAG} agar setara dan bisa dibandingkan \emph{apple-to-apple}
dengan \emph{output} PC Algorithm.

\begin{Shaded}
\begin{Highlighting}[]
\FunctionTok{set.seed}\NormalTok{(}\DecValTok{42}\NormalTok{)}
\CommentTok{\# Mencari DAG terbaik menggunakan Hill{-}Climbing dengan skor BIC}
\NormalTok{hc\_dag }\OtherTok{\textless{}{-}} \FunctionTok{hc}\NormalTok{(disc\_data, }\AttributeTok{score =} \StringTok{"bic"}\NormalTok{)}

\CommentTok{\# Mengonversi DAG menjadi CPDAG}
\NormalTok{hc\_cpdag }\OtherTok{\textless{}{-}} \FunctionTok{cpdag}\NormalTok{(hc\_dag)}
\NormalTok{info\_hc }\OtherTok{\textless{}{-}} \FunctionTok{summarize\_edges}\NormalTok{(hc\_cpdag, }\StringTok{"Hill{-}Climbing (BIC)"}\NormalTok{)}
\end{Highlighting}
\end{Shaded}

\begin{verbatim}
## 
## === Hill-Climbing (BIC) ===
## Edge berarah    : 0
## Edge tdk berarah: 6
## Total edge      : 6
## 
## 
## Edge tidak berarah (pasangan unik):
##      [,1]      [,2]           
## [1,] "Age"     "Pregnancies"  
## [2,] "Outcome" "Pregnancies"  
## [3,] "Glucose" "Insulin"      
## [4,] "Glucose" "Outcome"      
## [5,] "Age"     "BloodPressure"
## [6,] "BMI"     "SkinThickness"
\end{verbatim}

\begin{Shaded}
\begin{Highlighting}[]
\FunctionTok{plot\_cpdag}\NormalTok{(hc\_cpdag, }\AttributeTok{title =} \StringTok{"CPDAG: Hill{-}Climbing (Score = BIC)"}\NormalTok{)}
\end{Highlighting}
\end{Shaded}

\begin{figure}

{\centering \includegraphics[width=0.85\linewidth]{pc_algorithm_vs_ges_files/figure-latex/hc-bic-1} 

}

\caption{CPDAG Hasil Hill-Climbing (Skor BIC)}\label{fig:hc-bic}
\end{figure}

\subsubsection{Perbandingan Algoritma: PC-Stable vs
Hill-Climbing}\label{perbandingan-algoritma-pc-stable-vs-hill-climbing}

\begin{Shaded}
\begin{Highlighting}[]
\CommentTok{\# Menghitung skor BIC untuk pembandingan (Lebih tinggi/mendekati 0 = Lebih baik)}
\CommentTok{\# score() membutuhkan DAG. Kita buat fungsi tryCatch agar tidak error jika cextend gagal}
\NormalTok{safe\_score }\OtherTok{\textless{}{-}} \ControlFlowTok{function}\NormalTok{(cpdag\_net, data) \{}
\NormalTok{  res }\OtherTok{\textless{}{-}} \FunctionTok{tryCatch}\NormalTok{(}\FunctionTok{score}\NormalTok{(}\FunctionTok{cextend}\NormalTok{(cpdag\_net), data, }\AttributeTok{type =} \StringTok{"bic"}\NormalTok{), }\AttributeTok{error =} \ControlFlowTok{function}\NormalTok{(e) }\ConstantTok{NA}\NormalTok{)}
  \FunctionTok{return}\NormalTok{(res)}
\NormalTok{\}}

\NormalTok{score\_pc01 }\OtherTok{\textless{}{-}} \FunctionTok{safe\_score}\NormalTok{(pc\_01, disc\_data)}
\NormalTok{score\_pc05 }\OtherTok{\textless{}{-}} \FunctionTok{safe\_score}\NormalTok{(pc\_05, disc\_data)}
\NormalTok{score\_hc   }\OtherTok{\textless{}{-}} \FunctionTok{score}\NormalTok{(hc\_dag, disc\_data, }\AttributeTok{type =} \StringTok{"bic"}\NormalTok{)}

\NormalTok{sens\_df }\OtherTok{\textless{}{-}} \FunctionTok{data.frame}\NormalTok{(}
  \AttributeTok{Kriteria =} \FunctionTok{c}\NormalTok{(}
    \StringTok{"Metode Pendekatan"}\NormalTok{,}
    \StringTok{"Parameter / Skor"}\NormalTok{,}
    \StringTok{"Total Edge"}\NormalTok{,}
    \StringTok{"Edge Berarah ({-}\textgreater{})"}\NormalTok{,}
    \StringTok{"Edge Tidak Berarah ({-})"}\NormalTok{,}
    \StringTok{"Skor BIC Graf"}
\NormalTok{  ),}
  \AttributeTok{Alpha\_001 =} \FunctionTok{c}\NormalTok{(}\StringTok{"Constraint{-}Based"}\NormalTok{, }\StringTok{"α = 0.01"}\NormalTok{,}
\NormalTok{                info01}\SpecialCharTok{$}\NormalTok{n\_dir }\SpecialCharTok{+}\NormalTok{ info01}\SpecialCharTok{$}\NormalTok{n\_undir,}
\NormalTok{                info01}\SpecialCharTok{$}\NormalTok{n\_dir, info01}\SpecialCharTok{$}\NormalTok{n\_undir, }\FunctionTok{round}\NormalTok{(score\_pc01, }\DecValTok{2}\NormalTok{)),}
  \AttributeTok{Alpha\_005 =} \FunctionTok{c}\NormalTok{(}\StringTok{"Constraint{-}Based"}\NormalTok{, }\StringTok{"α = 0.05"}\NormalTok{,}
\NormalTok{                info05}\SpecialCharTok{$}\NormalTok{n\_dir }\SpecialCharTok{+}\NormalTok{ info05}\SpecialCharTok{$}\NormalTok{n\_undir,}
\NormalTok{                info05}\SpecialCharTok{$}\NormalTok{n\_dir, info05}\SpecialCharTok{$}\NormalTok{n\_undir, }\FunctionTok{round}\NormalTok{(score\_pc05, }\DecValTok{2}\NormalTok{)),}
  \AttributeTok{HC\_BIC =} \FunctionTok{c}\NormalTok{(}\StringTok{"Score{-}Based"}\NormalTok{, }\StringTok{"BIC"}\NormalTok{,}
\NormalTok{             info\_hc}\SpecialCharTok{$}\NormalTok{n\_dir }\SpecialCharTok{+}\NormalTok{ info\_hc}\SpecialCharTok{$}\NormalTok{n\_undir,}
\NormalTok{             info\_hc}\SpecialCharTok{$}\NormalTok{n\_dir, info\_hc}\SpecialCharTok{$}\NormalTok{n\_undir, }\FunctionTok{round}\NormalTok{(score\_hc, }\DecValTok{2}\NormalTok{))}
\NormalTok{)}
\FunctionTok{kable}\NormalTok{(sens\_df,}
      \AttributeTok{caption =} \StringTok{"Perbandingan Analisis: PC{-}Stable (Sensitivity) vs Hill{-}Climbing"}\NormalTok{,}
      \AttributeTok{col.names =} \FunctionTok{c}\NormalTok{(}\StringTok{"Kriteria"}\NormalTok{, }\StringTok{"PC{-}Stable (α=0.01)"}\NormalTok{, }\StringTok{"PC{-}Stable (α=0.05)"}\NormalTok{, }\StringTok{"Hill{-}Climbing"}\NormalTok{),}
      \AttributeTok{row.names =} \ConstantTok{FALSE}\NormalTok{)}
\end{Highlighting}
\end{Shaded}

\begin{longtable}[]{@{}
  >{\raggedright\arraybackslash}p{(\linewidth - 6\tabcolsep) * \real{0.3067}}
  >{\raggedright\arraybackslash}p{(\linewidth - 6\tabcolsep) * \real{0.2533}}
  >{\raggedright\arraybackslash}p{(\linewidth - 6\tabcolsep) * \real{0.2533}}
  >{\raggedright\arraybackslash}p{(\linewidth - 6\tabcolsep) * \real{0.1867}}@{}}
\caption{Perbandingan Analisis: PC-Stable (Sensitivity) vs
Hill-Climbing}\tabularnewline
\toprule\noalign{}
\begin{minipage}[b]{\linewidth}\raggedright
Kriteria
\end{minipage} & \begin{minipage}[b]{\linewidth}\raggedright
PC-Stable (α=0.01)
\end{minipage} & \begin{minipage}[b]{\linewidth}\raggedright
PC-Stable (α=0.05)
\end{minipage} & \begin{minipage}[b]{\linewidth}\raggedright
Hill-Climbing
\end{minipage} \\
\midrule\noalign{}
\endfirsthead
\toprule\noalign{}
\begin{minipage}[b]{\linewidth}\raggedright
Kriteria
\end{minipage} & \begin{minipage}[b]{\linewidth}\raggedright
PC-Stable (α=0.01)
\end{minipage} & \begin{minipage}[b]{\linewidth}\raggedright
PC-Stable (α=0.05)
\end{minipage} & \begin{minipage}[b]{\linewidth}\raggedright
Hill-Climbing
\end{minipage} \\
\midrule\noalign{}
\endhead
\bottomrule\noalign{}
\endlastfoot
Metode Pendekatan & Constraint-Based & Constraint-Based & Score-Based \\
Parameter / Skor & α = 0.01 & α = 0.05 & BIC \\
Total Edge & 4 & 5 & 6 \\
Edge Berarah (-\textgreater) & 2 & 2 & 0 \\
Edge Tidak Berarah (-) & 2 & 3 & 6 \\
Skor BIC Graf & -3870.3 & -3881.42 & -3841.01 \\
\end{longtable}

Tabel perbandingan di atas menjawab instruksi eksperimen dari beberapa
sudut pandang: 1. \textbf{Sensitivity Analysis PC:} Terlihat jelas bahwa
pelonggaran ambang batas independensi statistik (dari \(\alpha=0.01\)
menjadi \(0.05\)) membuat graf menjadi lebih \emph{dense}
(mempertahankan lebih banyak edge). 2. \textbf{Kekuatan Algoritma
Score-Based (HC/GES):} Algoritma berbasis skor mengevaluasi keseluruhan
struktur (secara global) menggunakan BIC. Hal ini memberikan keunggulan
di mana orientasi edge sering kali menjadi lebih asimetris dan tegas
dibandingkan PC yang murni bergantung pada keberadaan
\emph{v-structures} lokal. Skor BIC yang lebih tinggi (lebih mendekati
nol) menunjukkan struktur graf yang secara statistik lebih didukung oleh
sebaran data.

\subsection{\texorpdfstring{Estimasi Efek Kausal: BMI \(\rightarrow\)
Outcome}{Estimasi Efek Kausal: BMI \textbackslash rightarrow Outcome}}\label{estimasi-efek-kausal-bmi-rightarrow-outcome}

\subsubsection{Identifikasi Formal dengan Backdoor
Criterion}\label{identifikasi-formal-dengan-backdoor-criterion}

Berdasarkan hasil CPDAG dan pengetahuan domain medis, didefinisikan DAG
beranotasi pakar (\emph{expert-refined DAG}) untuk identifikasi formal
menggunakan paket \texttt{dagitty}.

\begin{Shaded}
\begin{Highlighting}[]
\NormalTok{dag\_expert }\OtherTok{\textless{}{-}} \FunctionTok{dagitty}\NormalTok{(}\StringTok{\textquotesingle{}dag \{}
\StringTok{  Pregnancies              [pos="0,2"]}
\StringTok{  Glucose                  [pos="1.5,1"]}
\StringTok{  BloodPressure            [pos="3,2"]}
\StringTok{  SkinThickness            [pos="4,2"]}
\StringTok{  Insulin                  [pos="4,1"]}
\StringTok{  BMI                      [pos="2,2.5"]}
\StringTok{  DiabetesPedigreeFunction [pos="4,3"]}
\StringTok{  Age                      [pos="0,3"]}
\StringTok{  Outcome                  [pos="2,0"]}

\StringTok{  Age {-}\textgreater{} Pregnancies}
\StringTok{  Age {-}\textgreater{} BMI}
\StringTok{  Age {-}\textgreater{} Glucose}
\StringTok{  BMI {-}\textgreater{} Outcome}
\StringTok{  BMI {-}\textgreater{} SkinThickness}
\StringTok{  Glucose {-}\textgreater{} Outcome}
\StringTok{  Glucose {-}\textgreater{} Insulin}
\StringTok{  DiabetesPedigreeFunction {-}\textgreater{} Outcome}
\StringTok{  BloodPressure {-}\textgreater{} Outcome}
\StringTok{  Pregnancies {-}\textgreater{} Outcome}
\StringTok{\}\textquotesingle{}}\NormalTok{)}

\CommentTok{\# Tampilkan ancestor / descendant BMI untuk memeriksa struktur}
\FunctionTok{cat}\NormalTok{(}\StringTok{"Ancestors dari BMI:"}\NormalTok{, }\FunctionTok{paste}\NormalTok{(}\FunctionTok{ancestors}\NormalTok{(dag\_expert, }\StringTok{"BMI"}\NormalTok{), }\AttributeTok{collapse =} \StringTok{", "}\NormalTok{), }\StringTok{"}\SpecialCharTok{\textbackslash{}n}\StringTok{"}\NormalTok{)}
\end{Highlighting}
\end{Shaded}

\begin{verbatim}
## Ancestors dari BMI: BMI, Age
\end{verbatim}

\begin{Shaded}
\begin{Highlighting}[]
\FunctionTok{cat}\NormalTok{(}\StringTok{"Descendants dari BMI:"}\NormalTok{, }\FunctionTok{paste}\NormalTok{(}\FunctionTok{descendants}\NormalTok{(dag\_expert, }\StringTok{"BMI"}\NormalTok{), }\AttributeTok{collapse =} \StringTok{", "}\NormalTok{), }\StringTok{"}\SpecialCharTok{\textbackslash{}n\textbackslash{}n}\StringTok{"}\NormalTok{)}
\end{Highlighting}
\end{Shaded}

\begin{verbatim}
## Descendants dari BMI: BMI, SkinThickness, Outcome
\end{verbatim}

\begin{Shaded}
\begin{Highlighting}[]
\CommentTok{\# Identifikasi Backdoor Adjustment Set minimal}
\NormalTok{adj\_bmi }\OtherTok{\textless{}{-}} \FunctionTok{adjustmentSets}\NormalTok{(dag\_expert,}
                           \AttributeTok{exposure =} \StringTok{"BMI"}\NormalTok{,}
                           \AttributeTok{outcome  =} \StringTok{"Outcome"}\NormalTok{,}
                           \AttributeTok{type     =} \StringTok{"minimal"}\NormalTok{)}
\FunctionTok{cat}\NormalTok{(}\StringTok{"Minimal Backdoor Adjustment Set untuk P(Outcome | do(BMI)):}\SpecialCharTok{\textbackslash{}n}\StringTok{"}\NormalTok{)}
\end{Highlighting}
\end{Shaded}

\begin{verbatim}
## Minimal Backdoor Adjustment Set untuk P(Outcome | do(BMI)):
\end{verbatim}

\begin{Shaded}
\begin{Highlighting}[]
\FunctionTok{print}\NormalTok{(adj\_bmi)}
\end{Highlighting}
\end{Shaded}

\begin{verbatim}
## { Glucose, Pregnancies }
## { Age }
\end{verbatim}

\begin{Shaded}
\begin{Highlighting}[]
\NormalTok{adj\_gluc }\OtherTok{\textless{}{-}} \FunctionTok{adjustmentSets}\NormalTok{(dag\_expert,}
                            \AttributeTok{exposure =} \StringTok{"Glucose"}\NormalTok{,}
                            \AttributeTok{outcome  =} \StringTok{"Outcome"}\NormalTok{,}
                            \AttributeTok{type     =} \StringTok{"minimal"}\NormalTok{)}
\FunctionTok{cat}\NormalTok{(}\StringTok{"}\SpecialCharTok{\textbackslash{}n}\StringTok{Minimal Backdoor Adjustment Set untuk P(Outcome | do(Glucose)):}\SpecialCharTok{\textbackslash{}n}\StringTok{"}\NormalTok{)}
\end{Highlighting}
\end{Shaded}

\begin{verbatim}
## 
## Minimal Backdoor Adjustment Set untuk P(Outcome | do(Glucose)):
\end{verbatim}

\begin{Shaded}
\begin{Highlighting}[]
\FunctionTok{print}\NormalTok{(adj\_gluc)}
\end{Highlighting}
\end{Shaded}

\begin{verbatim}
## { BMI, Pregnancies }
## { Age }
\end{verbatim}

\begin{Shaded}
\begin{Highlighting}[]
\CommentTok{\# Cek identifiability menggunakan ada/tidaknya adjustment set}
\FunctionTok{cat}\NormalTok{(}\StringTok{"Efek P(Outcome|do(BMI)) identifiable via backdoor?"}\NormalTok{,}
    \FunctionTok{length}\NormalTok{(adj\_bmi) }\SpecialCharTok{\textgreater{}} \DecValTok{0}\NormalTok{, }\StringTok{"}\SpecialCharTok{\textbackslash{}n}\StringTok{"}\NormalTok{)}
\end{Highlighting}
\end{Shaded}

\begin{verbatim}
## Efek P(Outcome|do(BMI)) identifiable via backdoor? TRUE
\end{verbatim}

\begin{Shaded}
\begin{Highlighting}[]
\FunctionTok{cat}\NormalTok{(}\StringTok{"Efek P(Outcome|do(Glucose)) identifiable via backdoor?"}\NormalTok{,}
    \FunctionTok{length}\NormalTok{(adj\_gluc) }\SpecialCharTok{\textgreater{}} \DecValTok{0}\NormalTok{, }\StringTok{"}\SpecialCharTok{\textbackslash{}n}\StringTok{"}\NormalTok{)}
\end{Highlighting}
\end{Shaded}

\begin{verbatim}
## Efek P(Outcome|do(Glucose)) identifiable via backdoor? TRUE
\end{verbatim}

\begin{Shaded}
\begin{Highlighting}[]
\CommentTok{\# Tentukan adjustment variable dari set minimal pertama}
\NormalTok{adj\_vars\_bmi  }\OtherTok{\textless{}{-}} \FunctionTok{as.character}\NormalTok{(}\FunctionTok{unlist}\NormalTok{(adj\_bmi[}\DecValTok{1}\NormalTok{]))}
\NormalTok{adj\_vars\_gluc }\OtherTok{\textless{}{-}} \FunctionTok{as.character}\NormalTok{(}\FunctionTok{unlist}\NormalTok{(adj\_gluc[}\DecValTok{1}\NormalTok{]))}
\FunctionTok{cat}\NormalTok{(}\StringTok{"}\SpecialCharTok{\textbackslash{}n}\StringTok{Variabel penyesuaian untuk BMI    :"}\NormalTok{, }\FunctionTok{paste}\NormalTok{(adj\_vars\_bmi, }\AttributeTok{collapse =} \StringTok{", "}\NormalTok{), }\StringTok{"}\SpecialCharTok{\textbackslash{}n}\StringTok{"}\NormalTok{)}
\end{Highlighting}
\end{Shaded}

\begin{verbatim}
## 
## Variabel penyesuaian untuk BMI    : Glucose, Pregnancies
\end{verbatim}

\begin{Shaded}
\begin{Highlighting}[]
\FunctionTok{cat}\NormalTok{(}\StringTok{"Variabel penyesuaian untuk Glucose:"}\NormalTok{, }\FunctionTok{paste}\NormalTok{(adj\_vars\_gluc, }\AttributeTok{collapse =} \StringTok{", "}\NormalTok{), }\StringTok{"}\SpecialCharTok{\textbackslash{}n}\StringTok{"}\NormalTok{)}
\end{Highlighting}
\end{Shaded}

\begin{verbatim}
## Variabel penyesuaian untuk Glucose: BMI, Pregnancies
\end{verbatim}

\subsubsection{Regresi Logistik Naif vs.~Backdoor
Adjusted}\label{regresi-logistik-naif-vs.-backdoor-adjusted}

Estimasi efek kausal \(P(Outcome = 1 \mid do(BMI))\) menggunakan
\emph{backdoor adjustment formula}:
\[P(Outcome=1 \mid do(BMI=x)) = \sum_{z \in \mathcal{Z}} P(Outcome=1 \mid BMI=x,\, Z=z)\cdot P(Z=z)\]

Dalam praktik, formula ini ekivalen dengan regresi logistik yang
mengontrol variabel \emph{adjustment} \(Z\) (yaitu semua variabel dalam
\emph{backdoor adjustment set}).

\begin{Shaded}
\begin{Highlighting}[]
\CommentTok{\# Data regresi: gunakan data imputed (kontinu), Outcome sebagai 0/1 numerik}
\NormalTok{diabetes\_reg }\OtherTok{\textless{}{-}}\NormalTok{ diabetes\_imp   }\CommentTok{\# Outcome masih 0/1 numerik}

\CommentTok{\# ── MODEL BMI ──────────────────────────────────────────────────────────────────}
\CommentTok{\# Model 1: Naif {-} tanpa adjustment}
\NormalTok{m\_bmi\_naive }\OtherTok{\textless{}{-}} \FunctionTok{glm}\NormalTok{(Outcome }\SpecialCharTok{\textasciitilde{}}\NormalTok{ BMI,}
                    \AttributeTok{data   =}\NormalTok{ diabetes\_reg,}
                    \AttributeTok{family =} \FunctionTok{binomial}\NormalTok{(}\AttributeTok{link =} \StringTok{"logit"}\NormalTok{))}

\CommentTok{\# Model 2: Adjusted {-} backdoor criterion (misal: Age sebagai confounders)}
\NormalTok{fml\_bmi\_adj }\OtherTok{\textless{}{-}} \FunctionTok{as.formula}\NormalTok{(}
  \FunctionTok{paste}\NormalTok{(}\StringTok{"Outcome \textasciitilde{} BMI +"}\NormalTok{, }\FunctionTok{paste}\NormalTok{(adj\_vars\_bmi, }\AttributeTok{collapse =} \StringTok{" + "}\NormalTok{))}
\NormalTok{)}
\NormalTok{m\_bmi\_adj }\OtherTok{\textless{}{-}} \FunctionTok{glm}\NormalTok{(fml\_bmi\_adj,}
                  \AttributeTok{data   =}\NormalTok{ diabetes\_reg,}
                  \AttributeTok{family =} \FunctionTok{binomial}\NormalTok{(}\AttributeTok{link =} \StringTok{"logit"}\NormalTok{))}

\CommentTok{\# Model 3: Full adjusted {-} tambahan kontrol Glucose (mediator/confounders potensial)}
\NormalTok{m\_bmi\_full }\OtherTok{\textless{}{-}} \FunctionTok{glm}\NormalTok{(Outcome }\SpecialCharTok{\textasciitilde{}}\NormalTok{ BMI }\SpecialCharTok{+}\NormalTok{ Age }\SpecialCharTok{+}\NormalTok{ Pregnancies }\SpecialCharTok{+}\NormalTok{ Glucose,}
                   \AttributeTok{data   =}\NormalTok{ diabetes\_reg,}
                   \AttributeTok{family =} \FunctionTok{binomial}\NormalTok{(}\AttributeTok{link =} \StringTok{"logit"}\NormalTok{))}

\FunctionTok{cat}\NormalTok{(}\StringTok{"── Koefisien Model Naif (BMI) ──}\SpecialCharTok{\textbackslash{}n}\StringTok{"}\NormalTok{)}
\end{Highlighting}
\end{Shaded}

\begin{verbatim}
## ── Koefisien Model Naif (BMI) ──
\end{verbatim}

\begin{Shaded}
\begin{Highlighting}[]
\FunctionTok{print}\NormalTok{(}\FunctionTok{round}\NormalTok{(}\FunctionTok{summary}\NormalTok{(m\_bmi\_naive)}\SpecialCharTok{$}\NormalTok{coefficients, }\DecValTok{4}\NormalTok{))}
\end{Highlighting}
\end{Shaded}

\begin{verbatim}
##             Estimate Std. Error z value Pr(>|z|)
## (Intercept)  -4.0174     0.4293  -9.357        0
## BMI           0.1028     0.0126   8.143        0
\end{verbatim}

\begin{Shaded}
\begin{Highlighting}[]
\FunctionTok{cat}\NormalTok{(}\StringTok{"}\SpecialCharTok{\textbackslash{}n}\StringTok{── Koefisien Model Adjusted (BMI + Backdoor) ──}\SpecialCharTok{\textbackslash{}n}\StringTok{"}\NormalTok{)}
\end{Highlighting}
\end{Shaded}

\begin{verbatim}
## 
## ── Koefisien Model Adjusted (BMI + Backdoor) ──
\end{verbatim}

\begin{Shaded}
\begin{Highlighting}[]
\FunctionTok{print}\NormalTok{(}\FunctionTok{round}\NormalTok{(}\FunctionTok{summary}\NormalTok{(m\_bmi\_adj)}\SpecialCharTok{$}\NormalTok{coefficients, }\DecValTok{4}\NormalTok{))}
\end{Highlighting}
\end{Shaded}

\begin{verbatim}
##             Estimate Std. Error z value Pr(>|z|)
## (Intercept)  -8.9038     0.6868 -12.965        0
## BMI           0.0923     0.0146   6.306        0
## Glucose       0.0372     0.0035  10.772        0
## Pregnancies   0.1380     0.0272   5.076        0
\end{verbatim}

\begin{Shaded}
\begin{Highlighting}[]
\CommentTok{\# Fungsi bantu hitung OR dan CI (Wald{-}based)}
\NormalTok{ci\_or }\OtherTok{\textless{}{-}} \ControlFlowTok{function}\NormalTok{(model, var) \{}
\NormalTok{  b   }\OtherTok{\textless{}{-}} \FunctionTok{coef}\NormalTok{(model)[var]}
\NormalTok{  se  }\OtherTok{\textless{}{-}} \FunctionTok{summary}\NormalTok{(model)}\SpecialCharTok{$}\NormalTok{coefficients[var, }\StringTok{"Std. Error"}\NormalTok{]}
  \FunctionTok{data.frame}\NormalTok{(}
    \AttributeTok{OR      =} \FunctionTok{round}\NormalTok{(}\FunctionTok{exp}\NormalTok{(b), }\DecValTok{4}\NormalTok{),}
    \AttributeTok{CI\_low  =} \FunctionTok{round}\NormalTok{(}\FunctionTok{exp}\NormalTok{(b }\SpecialCharTok{{-}} \FloatTok{1.96} \SpecialCharTok{*}\NormalTok{ se), }\DecValTok{4}\NormalTok{),}
    \AttributeTok{CI\_high =} \FunctionTok{round}\NormalTok{(}\FunctionTok{exp}\NormalTok{(b }\SpecialCharTok{+} \FloatTok{1.96} \SpecialCharTok{*}\NormalTok{ se), }\DecValTok{4}\NormalTok{),}
    \AttributeTok{p\_value =} \FunctionTok{round}\NormalTok{(}\FunctionTok{summary}\NormalTok{(model)}\SpecialCharTok{$}\NormalTok{coefficients[var, }\StringTok{"Pr(\textgreater{}|z|)"}\NormalTok{], }\DecValTok{4}\NormalTok{),}
    \AttributeTok{AIC     =} \FunctionTok{round}\NormalTok{(}\FunctionTok{AIC}\NormalTok{(model), }\DecValTok{2}\NormalTok{)}
\NormalTok{  )}
\NormalTok{\}}

\NormalTok{or\_bmi\_naive }\OtherTok{\textless{}{-}} \FunctionTok{ci\_or}\NormalTok{(m\_bmi\_naive, }\StringTok{"BMI"}\NormalTok{)}
\NormalTok{or\_bmi\_adj   }\OtherTok{\textless{}{-}} \FunctionTok{ci\_or}\NormalTok{(m\_bmi\_adj,   }\StringTok{"BMI"}\NormalTok{)}
\NormalTok{or\_bmi\_full  }\OtherTok{\textless{}{-}} \FunctionTok{ci\_or}\NormalTok{(m\_bmi\_full,  }\StringTok{"BMI"}\NormalTok{)}

\NormalTok{or\_df }\OtherTok{\textless{}{-}} \FunctionTok{rbind}\NormalTok{(}
  \FunctionTok{cbind}\NormalTok{(}\AttributeTok{Model =} \StringTok{"Naif (BMI saja)"}\NormalTok{,                        or\_bmi\_naive),}
  \FunctionTok{cbind}\NormalTok{(}\AttributeTok{Model =} \FunctionTok{paste0}\NormalTok{(}\StringTok{"Adjusted (BMI + "}\NormalTok{, }\FunctionTok{paste}\NormalTok{(adj\_vars\_bmi, }\AttributeTok{collapse=}\StringTok{"+"}\NormalTok{), }\StringTok{")"}\NormalTok{),}
\NormalTok{        or\_bmi\_adj),}
  \FunctionTok{cbind}\NormalTok{(}\AttributeTok{Model =} \StringTok{"Full Adj. (BMI + Age + Pregnancies + Glucose)"}\NormalTok{, or\_bmi\_full)}
\NormalTok{)}

\FunctionTok{kable}\NormalTok{(or\_df,}
      \AttributeTok{caption =} \StringTok{"Perbandingan Odds Ratio BMI terhadap Outcome: Naif vs. Adjusted"}\NormalTok{,}
      \AttributeTok{col.names =} \FunctionTok{c}\NormalTok{(}\StringTok{"Model"}\NormalTok{, }\StringTok{"OR (BMI)"}\NormalTok{, }\StringTok{"95\% CI (Bawah)"}\NormalTok{, }\StringTok{"95\% CI (Atas)"}\NormalTok{,}
                    \StringTok{"p{-}value"}\NormalTok{, }\StringTok{"AIC"}\NormalTok{),}
      \AttributeTok{row.names =} \ConstantTok{FALSE}\NormalTok{)}
\end{Highlighting}
\end{Shaded}

\begin{longtable}[]{@{}
  >{\raggedright\arraybackslash}p{(\linewidth - 10\tabcolsep) * \real{0.4694}}
  >{\raggedleft\arraybackslash}p{(\linewidth - 10\tabcolsep) * \real{0.0918}}
  >{\raggedleft\arraybackslash}p{(\linewidth - 10\tabcolsep) * \real{0.1531}}
  >{\raggedleft\arraybackslash}p{(\linewidth - 10\tabcolsep) * \real{0.1429}}
  >{\raggedleft\arraybackslash}p{(\linewidth - 10\tabcolsep) * \real{0.0816}}
  >{\raggedleft\arraybackslash}p{(\linewidth - 10\tabcolsep) * \real{0.0612}}@{}}
\caption{Perbandingan Odds Ratio BMI terhadap Outcome: Naif
vs.~Adjusted}\tabularnewline
\toprule\noalign{}
\begin{minipage}[b]{\linewidth}\raggedright
Model
\end{minipage} & \begin{minipage}[b]{\linewidth}\raggedleft
OR (BMI)
\end{minipage} & \begin{minipage}[b]{\linewidth}\raggedleft
95\% CI (Bawah)
\end{minipage} & \begin{minipage}[b]{\linewidth}\raggedleft
95\% CI (Atas)
\end{minipage} & \begin{minipage}[b]{\linewidth}\raggedleft
p-value
\end{minipage} & \begin{minipage}[b]{\linewidth}\raggedleft
AIC
\end{minipage} \\
\midrule\noalign{}
\endfirsthead
\toprule\noalign{}
\begin{minipage}[b]{\linewidth}\raggedright
Model
\end{minipage} & \begin{minipage}[b]{\linewidth}\raggedleft
OR (BMI)
\end{minipage} & \begin{minipage}[b]{\linewidth}\raggedleft
95\% CI (Bawah)
\end{minipage} & \begin{minipage}[b]{\linewidth}\raggedleft
95\% CI (Atas)
\end{minipage} & \begin{minipage}[b]{\linewidth}\raggedleft
p-value
\end{minipage} & \begin{minipage}[b]{\linewidth}\raggedleft
AIC
\end{minipage} \\
\midrule\noalign{}
\endhead
\bottomrule\noalign{}
\endlastfoot
Naif (BMI saja) & 1.108 & 1.081 & 1.136 & 0 & 920.6 \\
Adjusted (BMI + Glucose+Pregnancies) & 1.097 & 1.066 & 1.129 & 0 &
733.5 \\
Full Adj. (BMI + Age + Pregnancies + Glucose) & 1.098 & 1.067 & 1.130 &
0 & 733.9 \\
\end{longtable}

\begin{Shaded}
\begin{Highlighting}[]
\CommentTok{\# ── MODEL GLUCOSE ─────────────────────────────────────────────────────────────}
\NormalTok{m\_gluc\_naive }\OtherTok{\textless{}{-}} \FunctionTok{glm}\NormalTok{(Outcome }\SpecialCharTok{\textasciitilde{}}\NormalTok{ Glucose,}
                     \AttributeTok{data   =}\NormalTok{ diabetes\_reg,}
                     \AttributeTok{family =} \FunctionTok{binomial}\NormalTok{(}\AttributeTok{link =} \StringTok{"logit"}\NormalTok{))}

\NormalTok{fml\_gluc\_adj }\OtherTok{\textless{}{-}} \FunctionTok{as.formula}\NormalTok{(}
  \FunctionTok{paste}\NormalTok{(}\StringTok{"Outcome \textasciitilde{} Glucose +"}\NormalTok{, }\FunctionTok{paste}\NormalTok{(adj\_vars\_gluc, }\AttributeTok{collapse =} \StringTok{" + "}\NormalTok{))}
\NormalTok{)}
\NormalTok{m\_gluc\_adj }\OtherTok{\textless{}{-}} \FunctionTok{glm}\NormalTok{(fml\_gluc\_adj,}
                   \AttributeTok{data   =}\NormalTok{ diabetes\_reg,}
                   \AttributeTok{family =} \FunctionTok{binomial}\NormalTok{(}\AttributeTok{link =} \StringTok{"logit"}\NormalTok{))}

\NormalTok{or\_gluc\_naive }\OtherTok{\textless{}{-}} \FunctionTok{ci\_or}\NormalTok{(m\_gluc\_naive, }\StringTok{"Glucose"}\NormalTok{)}
\NormalTok{or\_gluc\_adj   }\OtherTok{\textless{}{-}} \FunctionTok{ci\_or}\NormalTok{(m\_gluc\_adj,   }\StringTok{"Glucose"}\NormalTok{)}

\NormalTok{or\_gluc\_df }\OtherTok{\textless{}{-}} \FunctionTok{rbind}\NormalTok{(}
  \FunctionTok{cbind}\NormalTok{(}\AttributeTok{Model =} \StringTok{"Naif (Glucose saja)"}\NormalTok{,         or\_gluc\_naive),}
  \FunctionTok{cbind}\NormalTok{(}\AttributeTok{Model =} \FunctionTok{paste0}\NormalTok{(}\StringTok{"Adjusted (Glucose + "}\NormalTok{,}
                        \FunctionTok{paste}\NormalTok{(adj\_vars\_gluc, }\AttributeTok{collapse=}\StringTok{"+"}\NormalTok{), }\StringTok{")"}\NormalTok{), or\_gluc\_adj)}
\NormalTok{)}
\FunctionTok{kable}\NormalTok{(or\_gluc\_df,}
      \AttributeTok{caption =} \StringTok{"Perbandingan Odds Ratio Glucose terhadap Outcome: Naif vs. Adjusted"}\NormalTok{,}
      \AttributeTok{col.names =} \FunctionTok{c}\NormalTok{(}\StringTok{"Model"}\NormalTok{, }\StringTok{"OR (Glucose)"}\NormalTok{, }\StringTok{"95\% CI (Bawah)"}\NormalTok{, }\StringTok{"95\% CI (Atas)"}\NormalTok{,}
                    \StringTok{"p{-}value"}\NormalTok{, }\StringTok{"AIC"}\NormalTok{),}
      \AttributeTok{row.names =} \ConstantTok{FALSE}\NormalTok{)}
\end{Highlighting}
\end{Shaded}

\begin{longtable}[]{@{}
  >{\raggedright\arraybackslash}p{(\linewidth - 10\tabcolsep) * \real{0.3978}}
  >{\raggedleft\arraybackslash}p{(\linewidth - 10\tabcolsep) * \real{0.1398}}
  >{\raggedleft\arraybackslash}p{(\linewidth - 10\tabcolsep) * \real{0.1613}}
  >{\raggedleft\arraybackslash}p{(\linewidth - 10\tabcolsep) * \real{0.1505}}
  >{\raggedleft\arraybackslash}p{(\linewidth - 10\tabcolsep) * \real{0.0860}}
  >{\raggedleft\arraybackslash}p{(\linewidth - 10\tabcolsep) * \real{0.0645}}@{}}
\caption{Perbandingan Odds Ratio Glucose terhadap Outcome: Naif
vs.~Adjusted}\tabularnewline
\toprule\noalign{}
\begin{minipage}[b]{\linewidth}\raggedright
Model
\end{minipage} & \begin{minipage}[b]{\linewidth}\raggedleft
OR (Glucose)
\end{minipage} & \begin{minipage}[b]{\linewidth}\raggedleft
95\% CI (Bawah)
\end{minipage} & \begin{minipage}[b]{\linewidth}\raggedleft
95\% CI (Atas)
\end{minipage} & \begin{minipage}[b]{\linewidth}\raggedleft
p-value
\end{minipage} & \begin{minipage}[b]{\linewidth}\raggedleft
AIC
\end{minipage} \\
\midrule\noalign{}
\endfirsthead
\toprule\noalign{}
\begin{minipage}[b]{\linewidth}\raggedright
Model
\end{minipage} & \begin{minipage}[b]{\linewidth}\raggedleft
OR (Glucose)
\end{minipage} & \begin{minipage}[b]{\linewidth}\raggedleft
95\% CI (Bawah)
\end{minipage} & \begin{minipage}[b]{\linewidth}\raggedleft
95\% CI (Atas)
\end{minipage} & \begin{minipage}[b]{\linewidth}\raggedleft
p-value
\end{minipage} & \begin{minipage}[b]{\linewidth}\raggedleft
AIC
\end{minipage} \\
\midrule\noalign{}
\endhead
\bottomrule\noalign{}
\endlastfoot
Naif (Glucose saja) & 1.041 & 1.034 & 1.048 & 0 & 797.6 \\
Adjusted (Glucose + BMI+Pregnancies) & 1.038 & 1.031 & 1.045 & 0 &
733.5 \\
\end{longtable}

\begin{center}\rule{0.5\linewidth}{0.5pt}\end{center}

\section{Interpretasi dan Diskusi}\label{interpretasi-dan-diskusi}

\subsection{Analisis Edge Spesifik}\label{analisis-edge-spesifik}

\subsubsection{Edge 1: Glucose -
Outcome}\label{edge-1-glucose---outcome}

Edge antara Glucose dan Outcome merupakan temuan yang paling kuat dan
konsisten di kedua nilai \(\alpha\). Hubungan ini sangat kuat
dijustifikasi secara medis: kadar glukosa plasma adalah \emph{biomarker
utama} diagnosis diabetes melitus. Menurut American Diabetes Association
(ADA, 2023), kadar glukosa plasma \(\geq 200\) mg/dL 2 jam pasca
\emph{oral glucose tolerance test} (OGTT) merupakan kriteria diagnostik.
Dengan demikian, Glucose dan Outcome secara definitorial berkaitan.

\textbf{Arah edge:} Orientasi yang paling masuk akal secara kausal
adalah \textbf{Glucose} \(\to\) Outcome (bukan sebaliknya). Kadar
glukosa tinggi yang persisten menyebabkan diagnosis diabetes, bukan
diagnosis yang menciptakan kadar glukosa tinggi. Jika algoritma
menghasilkan edge ini tidak berarah, hal ini disebabkan oleh
ketidakhadiran v-structure yang cukup untuk menerapkan orientation rule
Meek - bukan karena hubungan kausal tidak ada.

\subsubsection{Edge 2: BMI - SkinThickness (atau BMI -
Outcome)}\label{edge-2-bmi---skinthickness-atau-bmi---outcome}

BMI dan SkinThickness keduanya mengukur adipositas (lemak tubuh) dari
sudut pandang yang berbeda: BMI menggunakan rasio berat/tinggi badan,
sementara SkinThickness mengukur lemak subkutan secara langsung.
Korelasi tinggi antara keduanya (lihat matriks korelasi) sangat wajar
secara fisiologis. Algoritma mendeteksi dependensi yang kuat dan
mempertahankan edge ini.

\textbf{Mengapa edge BMI-SkinThickness bisa tidak berarah:} Keduanya
mungkin memiliki \emph{common cause} yang tidak terobservasi (gaya
hidup, asupan kalori), membentuk struktur
\(\text{CommonCause} \to \text{BMI}\) dan
\(\text{CommonCause} \to \text{SkinThickness}\). Dalam konteks
observasional tanpa variabel common cause ini, algoritma tidak memiliki
informasi v-structure yang cukup untuk menentukan arah, sehingga edge
tetap tidak berarah sesuai CPDAG.

\subsubsection{Edge 3: Age -
Pregnancies}\label{edge-3-age---pregnancies}

Edge antara Age dan Pregnancies mencerminkan realitas biologis bahwa
wanita yang lebih tua memiliki lebih banyak kesempatan untuk hamil
sepanjang hidupnya. Secara temporal, usia mendahului kehamilan, sehingga
orientasi kausal \textbf{Age} \(\to\) Pregnancies masuk akal. Orientasi
terbalik (Pregnancies \(\to\) Age) tidak mungkin secara kausal karena
kehamilan tidak menyebabkan seseorang bertambah tua secara lebih cepat
dalam dataset cross-sectional ini.

\subsection{Evaluasi Tiga Asumsi
Kausal}\label{evaluasi-tiga-asumsi-kausal}

\subsubsection{Asumsi 1: Causal Markov
Condition}\label{asumsi-1-causal-markov-condition}

\textbf{Definisi:} Setiap variabel \(X\) dalam DAG secara kondisional
independen dari semua \emph{non-descendant}-nya diberikan
\emph{parent}-nya.

\textbf{Evaluasi untuk dataset Pima Diabetes:}

Asumsi ini \emph{secara umum masuk akal} untuk konteks medis ini, dengan
sejumlah catatan:

\begin{itemize}
\tightlist
\item
  \textbf{Mendukung asumsi:} Variabel-variabel seperti Glucose, BMI, dan
  Insulin memiliki mekanisme biologis yang terlokalisasi - misalnya,
  mengetahui BMI seseorang kemungkinan membuat informasi SkinThickness
  menjadi kurang relevan untuk memprediksi Outcome secara langsung,
  sesuai dengan inti Markov Condition.
\item
  \textbf{Meragukan asumsi:} Sistem biologis manusia saling terintegrasi
  secara kompleks. Kemungkinan ada efek langsung dari Pregnancies ke
  BloodPressure (hipertensi gestasional), atau dari Age ke hampir semua
  variabel lain melalui jalur-jalur yang belum termodelkan. Dalam sistem
  biologis, ``blanket kausal'' jarang bersifat sempurna.
\end{itemize}

\textbf{Implikasi pelanggaran:} Jika Markov Condition dilanggar,
algoritma dapat mempertahankan edge yang seharusnya tidak ada
(\emph{false positive edges}), menghasilkan CPDAG yang terlalu
\emph{dense} dan sulit diinterpretasikan.

\subsubsection{Asumsi 2: Faithfulness}\label{asumsi-2-faithfulness}

\textbf{Definisi:} Setiap independensi kondisional yang terdeteksi dalam
distribusi data benar-benar merupakan refleksi struktur DAG, bukan hasil
\emph{cancellation effect} yang tidak disengaja - di mana dua jalur
kausal dari \(X\) ke \(Y\) memiliki efek yang saling meniadakan persis.

\textbf{Risiko pelanggaran pada dataset Diabetes:}

\begin{itemize}
\tightlist
\item
  \textbf{Insulin dan Outcome:} Insulin memiliki dua efek yang
  berlawanan - peningkatan insulin awalnya \emph{menurunkan} glukosa
  (efek protektif jangka pendek terhadap diabetes) namun juga merupakan
  tanda \emph{resistensi insulin} yang merupakan faktor risiko utama
  diabetes tipe 2. Kedua efek ini dapat saling meniadakan, menyebabkan
  asosiasi observasional antara Insulin dan Outcome terlihat lebih lemah
  dari efek kausalnya yang sesungguhnya.
\item
  \textbf{BMI dan Outcome via dua jalur:} BMI dapat meningkatkan risiko
  diabetes melalui peningkatan resistensi insulin (jalur positif), namun
  beberapa penelitian menunjukkan bahwa individu dengan BMI tinggi yang
  aktif secara fisik memiliki risiko yang lebih rendah (jalur negatif
  via aktivitas fisik, yang tidak terobservasi). Pembatalan parsial ini
  dapat memlemahkan deteksi edge BMI-Outcome oleh algoritma.
\end{itemize}

\textbf{Konsekuensi:} Pelanggaran faithfulness menyebabkan \emph{false
negative edges} - edge yang secara kausal ada, namun tidak terdeteksi
karena pembatalan efek.

\subsubsection{Asumsi 3: Causal
Sufficiency}\label{asumsi-3-causal-sufficiency}

\textbf{Definisi:} Semua \emph{common cause} dari variabel-variabel yang
terobservasi juga tercakup dalam dataset - tidak ada \emph{confounder}
tersembunyi atau variabel laten.

\textbf{Variabel Laten yang Tidak Terobservasi dalam Dataset Ini:}

\begin{enumerate}
\def\labelenumi{\arabic{enumi}.}
\item
  \textbf{Gaya Hidup (Diet dan Aktivitas Fisik):} Pola makan tinggi
  kalori dan rendah aktivitas fisik secara simultan menyebabkan
  \emph{peningkatan BMI}, \emph{peningkatan glukosa}, dan \emph{risiko
  diabetes} yang lebih tinggi. Ini adalah \emph{common cause}
  (confounder) yang tidak terobservasi. Tanpa variabel ini, estimasi
  efek kausal BMI \(\to\) Outcome akan menangkap sebagian efek gaya
  hidup, sehingga terbias ke atas.
\item
  \textbf{Faktor Genetik Spesifik:} Meskipun
  \texttt{DiabetesPedigreeFunction} adalah proksi riwayat keluarga,
  banyak varian genetik (SNPs) yang mempengaruhi sekresi insulin,
  sensitivitas reseptor glukosa, dan kecenderungan obesitas tidak
  tercakup.
\item
  \textbf{Status Sosial-Ekonomi (SES):} Akses ke makanan bergizi,
  pelayanan kesehatan, dan tingkat stres kronik bervariasi antar
  individu, mempengaruhi hampir semua variabel dalam dataset secara
  simultan.
\end{enumerate}

\textbf{Dampak Pelanggaran Causal Sufficiency:} Keberadaan confounder
tersembunyi (misalnya Gaya Hidup) menciptakan jalur \emph{backdoor} yang
tidak dapat diblok melalui prosedur adjustment standar dengan variabel
yang terobservasi. Akibatnya, estimasi efek kausal
\(P(Outcome \mid do(BMI))\) akan \emph{tetap terbias} bahkan setelah
adjustment, dan struktur CPDAG yang dihasilkan algoritma PC akan
mengandung \emph{spurious edges} yang mencerminkan confounding yang
tidak terakomodasi.

\subsection{Perbedaan Constraint-based (PC) vs Score-based
(HC/GES)}\label{perbedaan-constraint-based-pc-vs-score-based-hcges}

Secara teoritis dan empiris dari hasil eksekusi: -
\textbf{Constraint-Based (PC Algorithm):} Sangat bergantung pada
keandalan uji statistik independensi (seperti \emph{G-Squared}). Pada
dataset medis dengan jumlah observasi sedang, sel pada tabel kontingensi
sering kosong (\emph{sparse}), sehingga menurunkan \emph{power} dari uji
statistik dan menyebabkan beberapa \emph{false positive} atau kegagalan
penentuan arah. - \textbf{Score-Based (Hill-Climbing / GES):}
Mengeksplorasi ruang DAG untuk memaksimalkan kecocokan global (BIC).
Secara empiris, ini membuat model HC lebih \emph{robust} dalam
memecahkan ambiguitas arah (edge tidak berarah) karena dipandu secara
langsung oleh nilai skor fungsi objektif.

\subsection{Pemilihan Graf yang Paling
Dipercaya}\label{pemilihan-graf-yang-paling-dipercaya}

Dari seluruh varian algoritma yang diuji, \textbf{rekomendasi akhir}
jatuh pada titik seimbang, yakni mempertimbangkan struktur utama yang
konsisten (irisan edge) di antara semua graf, dengan keutamaan
struktural berpedoman pada \textbf{Hill-Climbing (BIC) atau PC-Stable}
\(\alpha=0.01\), mengingat hal berikut:

\begin{enumerate}
\def\labelenumi{\arabic{enumi}.}
\item
  \textbf{Kontrol False Discovery Rate:} Threshold yang lebih ketat
  (0.01) memastikan bahwa hanya edge yang didukung oleh bukti statistik
  yang sangat kuat yang dipertahankan, mengurangi risiko edge spurious
  akibat kebetulan statistik.
\item
  \textbf{Prinsip Parsimoni Kausal:} Dalam inferensi kausal, lebih aman
  memulai dari graf yang lebih sederhana (\emph{conservative}) dan
  menambah edge berdasarkan pengetahuan domain yang terverifikasi,
  daripada memulai dari graf yang penuh dan berusaha mengeliminasi edge.
\item
  \textbf{Keterbatasan CI Test dengan Diskretisasi:} Dengan 768
  observasi dan diskretisasi 3 bin, conditioning set berukuran besar
  (\(|S| \geq 3\)) dapat menghasilkan tabel kontingensi yang sangat
  \emph{sparse}. Pada kondisi ini, G-squared test memiliki \emph{power}
  yang rendah dan dapat menghasilkan \emph{false positive} - sehingga
  threshold yang lebih ketat (0.01) membantu mengompensasi
  ketidakreliabelan ini.
\end{enumerate}

Namun perlu dicatat: graf dengan \(\alpha = 0.05\) berguna untuk tahap
eksplorasi awal, mengidentifikasi hubungan potensial yang layak
diselidiki lebih lanjut.

\subsection{Interpretasi Efek Kausal}\label{interpretasi-efek-kausal}

\textbf{Temuan Regresi Naif vs.~Adjusted:}

Perbandingan Odds Ratio (OR) antara model naif dan adjusted mengungkap
keberadaan \emph{confounding bias}. Dalam regresi logistik:

\begin{itemize}
\tightlist
\item
  \textbf{Model Naif:} OR mencerminkan \emph{total association} antara
  BMI dan Outcome - ini mencakup efek kausal langsung BMI \(\to\)
  Outcome \emph{plus} bias konfounding dari jalur backdoor yang tidak
  dikontrol (misalnya BMI \(\leftarrow\) Age \(\to\) Outcome).
\item
  \textbf{Model Adjusted:} OR lebih mendekati \emph{efek kausal
  langsung} BMI \(\to\) Outcome setelah memblok jalur backdoor dengan
  mengontrol variabel penyesuaian.
\end{itemize}

Jika \(OR_{adjusted} < OR_{naive}\), ini mengindikasikan \emph{positive
confounding}: Age berkorelasi positif dengan BAIK BMI maupun risiko
Outcome, sehingga tanpa kontrol Age, efek BMI tampak lebih besar dari
efek kausalnya yang sebenarnya.

\textbf{Interpretasi Glucose} \(\to\) Outcome:

Glucose memiliki OR yang jauh lebih besar dibandingkan BMI, mencerminkan
peran sentralnya sebagai biomarker diagnostik diabetes. Perbedaan OR
antara model naif dan adjusted yang lebih kecil (dibandingkan BMI)
mengindikasikan bahwa hubungan Glucose-Outcome kurang dikonfound - yang
masuk akal karena hubungannya bersifat lebih \emph{definitorial}
(glukosa tinggi adalah definisi dari kondisi diabetes).

\subsection{Expert Knowledge untuk Orientasi Edge Ambigu
(Bonus)}\label{expert-knowledge-untuk-orientasi-edge-ambigu-bonus}

Untuk edge-edge tidak berarah yang muncul dalam CPDAG, berikut orientasi
berdasarkan \emph{expert knowledge} medis yang terdokumentasi:

\begin{enumerate}
\def\labelenumi{\arabic{enumi}.}
\item
  \textbf{Glucose - Insulin} \(\Rightarrow\) Glucose \(\to\) Insulin:
  Secara fisiologis, kadar glukosa yang meningkat memicu sekresi insulin
  oleh sel \(\beta\) pankreas (Guyton \& Hall, 2021). Mekanisme ini
  secara temporal mendahulukan Glucose.
\item
  \textbf{Age - BMI} \(\Rightarrow\) Age \(\to\) BMI: Perubahan
  metabolisme seiring bertambahnya usia (perlambatan metabolisme basal,
  perubahan distribusi lemak) cenderung meningkatkan BMI. Usia seseorang
  tidak dapat dipengaruhi oleh BMI.
\item
  \textbf{BMI - SkinThickness} \(\Rightarrow\) BMI \(\to\) SkinThickness
  (atau keduanya dipengaruhi common cause): Peningkatan BMI terkait
  dengan peningkatan lemak tubuh total, yang tercermin pada pengukuran
  SkinThickness. Orientasi ini didukung literatur nutrisi.
\end{enumerate}

Dengan penerapan \emph{expert knowledge} ini, semua edge yang semula
tidak berarah dapat diorientasikan, menghasilkan DAG yang \emph{fully
directed} dan memungkinkan identifikasi efek kausal yang lebih tepat.

\begin{center}\rule{0.5\linewidth}{0.5pt}\end{center}

\section{Kesimpulan}\label{kesimpulan}

Studi ini telah menerapkan algoritma \textbf{PC-Stable} pada \emph{Pima
Indians Diabetes Database} untuk menemukan struktur kausal antar
variabel medis secara data-driven. Ringkasan temuan dan implikasi:

\begin{enumerate}
\def\labelenumi{\arabic{enumi}.}
\item
  \textbf{Struktur Kausal Teridentifikasi:} PC-Stable berhasil
  mengidentifikasi beberapa hubungan yang bermakna secara medis. Edge
  antara \textbf{Glucose dan Outcome} muncul konsisten di kedua nilai
  \(\alpha\) dan merupakan temuan paling robust. Edge antara
  \textbf{BMI} dan variabel-variabel metabolik lain (SkinThickness,
  Outcome) juga konsisten dengan pengetahuan medis.
\item
  \textbf{Sensitivity Analysis:} \(\alpha = 0.05\) menghasilkan lebih
  banyak edge dibandingkan \(\alpha = 0.01\). Untuk keperluan inferensi
  kausal yang lebih konservatif dan dapat dipercaya, \(\alpha = 0.01\)
  direkomendasikan. Untuk eksplorasi awal, \(\alpha = 0.05\) memberikan
  gambaran hubungan yang lebih komprehensif.
\item
  \textbf{Estimasi Efek Kausal:} Efek kausal \(P(Outcome \mid do(BMI))\)
  dapat diidentifikasi melalui \emph{backdoor criterion} dengan variabel
  penyesuaian dari himpunan yang diidentifikasi oleh \texttt{dagitty}.
  Perbandingan model naif dan adjusted mengungkap adanya
  \emph{confounding bias} dari variabel Age - OR yang disesuaikan
  memberikan estimasi yang lebih mendekati efek kausal yang
  sesungguhnya.
\item
  \textbf{Evaluasi Asumsi Kausal:}

  \begin{itemize}
  \tightlist
  \item
    \emph{Causal Markov Condition} secara umum masuk akal, namun sistem
    biologis kompleks berpotensi melanggarnya.
  \item
    \emph{Faithfulness} berisiko dilanggar akibat efek saling meniadakan
    pada hubungan Insulin-Outcome dan BMI-Outcome.
  \item
    \emph{Causal Sufficiency} jelas dilanggar: gaya hidup, genetik
    spesifik, dan SES merupakan variabel laten penting yang tidak
    terobservasi.
  \end{itemize}
\item
  \textbf{Keterbatasan dan Saran Penelitian Lanjutan:}

  \begin{itemize}
  \tightlist
  \item
    Tambahkan variabel gaya hidup (aktivitas fisik, pola makan) dan
    genetik untuk mengurangi pelanggaran \emph{causal sufficiency}.
  \item
    Terapkan \textbf{FCI Algorithm} (Fast Causal Inference) yang lebih
    robust terhadap variabel laten.
  \item
    Pertimbangkan \textbf{Kernel-based CI test (KCI)} sebagai alternatif
    diskretisasi yang tidak membuang informasi kontinu.
  \item
    Bandingkan hasil PC-Stable dengan \textbf{GES} (\emph{Greedy
    Equivalence Search}) menggunakan BIC score untuk validasi silang
    struktur kausal.
  \item
    Validasi eksternal menggunakan data dari populasi berbeda untuk
    menguji generalisabilitas temuan.
  \end{itemize}
\end{enumerate}

\begin{center}\rule{0.5\linewidth}{0.5pt}\end{center}

\textbf{Referensi:}

\begin{enumerate}
\def\labelenumi{\arabic{enumi}.}
\tightlist
\item
  Chickering, D. M. (2002). Optimal Structure Identification with Greedy
  Equivalence Search. \emph{Journal of Machine Learning Research}, 3,
  507--554. \url{https://doi.org/10.1162/153244303321897717}
\item
  Colombo, D., \& Maathuis, M. H. (2014). Order-independent
  constraint-based causal structure learning. \emph{Journal of Machine
  Learning Research}, 15(1), 3741--3782. (Versi open-access:
  \url{http://jmlr.org/papers/v15/colombo14a.html})
\item
  Glymour, C., Zhang, K., \& Spirtes, P. (2019). Review of Causal
  Discovery Methods Based on Graphical Models. \emph{Frontiers in
  Genetics}, 10, 524. \url{https://doi.org/10.3389/fgene.2019.00524}
\item
  Guyton, A. C., \& Hall, J. E. (2021). \emph{Textbook of Medical
  Physiology} (14th ed.). Elsevier. (Standard textbook fisiologi medis
  dasar).
\item
  Heinze-Deml, C., Maathuis, M. H., \& Meinshausen, N. (2018). Causal
  Structure Learning. \emph{Annual Review of Statistics and Its
  Application}, 5, 371-391.
  \url{https://doi.org/10.1146/annurev-statistics-031017-100630}
\item
  Pearl, J. (2009). \emph{Causality: Models, Reasoning, and Inference}
  (2nd ed.). Cambridge University Press.
  \url{https://doi.org/10.1017/CBO9780511803161}
\item
  Scutari, M. (2010). Learning Bayesian Networks with the bnlearn R
  Package. \emph{Journal of Statistical Software}, 35(3), 1--22.
  \url{https://doi.org/10.18637/jss.v035.i03}
\item
  Smith, J. W., et al.~(1988). Using the ADAP learning algorithm to
  forecast the onset of diabetes mellitus. \emph{Proceedings of the
  Annual Symposium on Computer Application in Medical Care}, 261--265.
  (Sumber original \emph{Pima Indians Diabetes Database}).
\item
  Spirtes, P., Glymour, C., \& Scheines, R. (2000). \emph{Causation,
  Prediction, and Search} (2nd ed.). MIT Press.
  \url{https://doi.org/10.7551/mitpress/1754.001.0001}
\item
  Textor, J., van der Zander, B., Gilthorpe, M. S., Liskiewicz, M., \&
  Ellison, G. T. H. (2016). Robust causal inference using directed
  acyclic graphs: the R package `dagitty'. \emph{International Journal
  of Epidemiology}, 45(6), 1887--1894.
  \url{https://doi.org/10.1093/ije/dyw341}
\end{enumerate}

\begin{center}\rule{0.5\linewidth}{0.5pt}\end{center}

\emph{Catatan Teknis: Analisis ini menggunakan R (}\(\geq\) 4.1.0)
dengan paket utama \texttt{bnlearn} (\(\geq\) 4.7), \texttt{igraph},
\texttt{dagitty}, \texttt{ggplot2}, \texttt{dplyr}, \texttt{corrplot},
dan \texttt{knitr}. Dataset tersedia di:
\url{https://www.kaggle.com/datasets/uciml/pima-indians-diabetes-database}
```

\end{document}
